% !TEX root = critical_path_ru.tex
\documentclass[12pt,a4paper]{article}
\usepackage{cmap}
\usepackage[margin=1in]{geometry}
\usepackage{amsmath,amssymb,amsfonts}
\usepackage[utf8]{inputenc}
\usepackage[T2A]{fontenc}
\usepackage[russian]{babel}
\usepackage[hidelinks]{hyperref}
\usepackage{comfortaa}
\renewcommand*\familydefault{\sfdefault}
\usepackage[x11names]{xcolor}

\makeatletter
\renewcommand{\thefootnote}{¹}
\renewcommand{\@makefnmark}{\mbox{\textcolor{white}{\thefootnote}}}
\renewcommand{\@makefntext}[1]{\parindent 1.8em\noindent\mbox{}\llap{\thefootnote}\textcolor{white}{#1}}
\renewcommand{\footnoterule}{\kern-3pt\color{white}\hrule width 0.4\textwidth height 0.4pt\kern 2.6pt}
\makeatother

\title{\textbf{Критический путь эволюции}}
\author{%
  Олег Понфиленок\\[0.35em]%
  {\normalsize Независимый исследователь; \href{mailto:ponfil@gmail.com}{ponfil@gmail.com}}%
}
\date{}

\begin{document}

\pagecolor{black}
\color{white}

\maketitle

\begin{abstract}
Эволюция жизни рассматривается как прохождение $N \sim 10^{14}$ малых высоковероятных шагов усложнения со средним временем шага $\mu$~--- критическим путём от теплового гейта реликтового излучения до технологической цивилизации за время $T \approx 13{,}8$~Gyr. При независимости времён шагов концентрация их суммы сжимает относительный разброс времён выхода как $1/\sqrt{N}$; в рамках модели отсюда следует дисперсия $\sigma = \sqrt{\mu T}$. Мы показываем, что наблюдаемое отсутствие «громких» цивилизаций (космическая тишина) даёт оценку межпланетной дисперсии времён выхода $\sigma \approx 800$--$1400$~лет~--- порядка горизонта перехода к детектируемости; инверсия тишины даёт $\mu \approx 25$--$75$~мин, что совпадает с генерационным временем микроорганизмов. Та же тишина ограничивает населённость Галактики~--- $n \sim 10^4$--$10^5$ цивилизационных планет, согласуясь с астрофизической оценкой по порядку величины. Два свойства критического пути~--- протекание при локальном избытке ресурсов и кинематическая уникальность самого быстрого маршрута~--- делают его параметры приближённо универсальными и обеспечивают синхронный темп развития жизни; космическая тишина указывает, что этот маршрут реализуется через углеродно-органическую химию жизни нашего типа. Сильнейшее следствие~--- конвергенция не только времён, но и 
формы цивилизации: технологические цивилизации на других планетах должны быть гуманоидами нашего функционального типа.
\end{abstract}

\noindent\textbf{Ключевые слова:} парадокс~Ферми; критический~путь~эволюции; модель~простых~шагов (против «трудных~шагов»); синхронизация~времён~появления~разума; MEPP; конвергентная~эволюция; населённость~Галактики; SETI

%%% ===================================================================
\section{Введение}
%%% ===================================================================

Парадокс Ферми~--- отсутствие наблюдаемых техносигнатур при ожидаемой обитаемости галактики~--- обычно обсуждается через редкие «трудные шаги» эволюции (модель Картера~\cite{Carter1983}) или через антропные аргументы о ранней позиции наблюдателя (модель grabby aliens~\cite{Hanson2021}). В данной работе предлагается иной, строго копернианский подход: эволюция жизни от рождения Вселенной до технологической цивилизации идёт по \emph{критическому пути}~--- лимитирующему, самому быстрому маршруту усложнения (по аналогии с методом критического пути в управлении проектами, critical-path method), а не по набору редких маловероятных переходов. Это \textbf{противоположно} понятию «critical steps» Картера: critical steps~--- мало и они трудны; critical path~--- много простых шагов, и важна лишь скорость лимитирующего маршрута.

Выбор в пользу простых шагов мотивирован не только умозрительно, но и буквальным прочтением Дарвиновского \emph{natura non facit saltum}~\cite{Darwin1859}: природа не делает скачков, усложнение~--- цепочка малых изменений, а не редких «трудных» переходов. Дополнительный стимул~--- асимметрия эмпирических данных: гипотеза «трудных шагов» сама нигде не доказана~--- неразложимость не была строго показана ни для одного предполагаемого кандидата,~--- тогда как при детальном изучении конкретные кандидаты (эукариогенез, происхождение жизни) раз за разом распадаются на цепочки кооперативных, высоковероятных подшагов~\cite{Mills2025}. Единственность происхождения сама по себе трудность не доказывает: возможны эффект инкумбента (первая успешная линия исключает повторные попытки), закрытие ниши средой и недонаблюдаемость вымерших альтернативных линий.

Термодинамически это можно понимать как отбор наиболее диссипативного маршрута: в открытой системе при избытке локальных ресурсов принцип максимального производства энтропии (MEPP~\cite{Martyushev,MartyushevSeleznev2006,Kleidon2010,Vallino2010}) выделяет путь, быстрее всего превращающий доступную свободную энергию в устойчивое усложнение. Поэтому критический путь~--- не только кинематически самый быстрый, но и физически наиболее диссипативный маршрут.

Эволюцию от теплового гейта реликтового излучения до технологической цивилизации мы рассматриваем как \emph{единую} цепочку $N$ очень многих простых шагов. При сложении многих независимых стадий среднее время растёт как $N$, а разброс~--- лишь как $\sqrt{N}$, поэтому \emph{относительный} разброс времён выхода сжимается с числом шагов: разные планеты приходят к цивилизации почти одновременно~--- в этом и состоит \textbf{синхронизация}. Она держится на двух свойствах критического пути: он \emph{единствен} и идёт с \emph{универсальным} тактом $\mu$ на плато скорости (при избытке ресурсов темп задаётся химией, а не средой). Количественный вывод этих свойств дан в \S~\ref{sec:steps}, порядковая оценка $\sigma$ из тишины~--- в \S~\ref{sec:inversion}; главная цель~--- показать качественный механизм синхронизации, а не точную калибровку временных масштабов.

По сути это иной взгляд на эволюцию~--- не как на череду случайностей, а как на закономерность: отдельный шаг случаен, но концентрация суммы огромного их числа делает итоговый путь практически детерминированным, и жизнь предстаёт не лотереей, а \emph{проектом} с единственным самым быстрым маршрутом усложнения. Поскольку этот маршрут кинематически уникален, модель предсказывает сильную конвергенцию не только времён выхода, но и \emph{формы цивилизации}: технологические цивилизации на других планетах~--- гуманоиды нашего функционального типа (\S~\ref{sec:time_form}).

Ближайший союзник по качественному тезису~--- работа Mills et al.~\cite{Mills2025} («нет трудных шагов»); вопрос о дисперсии времён появления разума ранее ставил Hair~\cite{Hair2011}, но с противоположным выводом: он подчёркивал \emph{большой} разброс и колоссальное временное преимущество древнейших цивилизаций, тогда как из суммы независимых шагов мы получаем узкий синхронный пик. Настоящая работа разделяет вывод «трудных шагов нет», но добавляет количественную рамку суммы последовательных стадий и оценку $\sigma$ из тишины (\S~\ref{sec:inversion}). От модели grabby aliens~\cite{Hanson2021} мы отличаемся отказом от степенного закона hard steps и от антропной посылки: синхронность выводится из концентрации суммы независимых шагов (аддитивность дисперсии) и локальной конвергентности космической фазы. Иную трактовку тишины дают Sandberg, Drexler и Ord~\cite{Sandberg2018}: парадокс «растворяется» статистически~--- при честном учёте многопорядковой неопределённости параметров уравнения Дрейка вероятность одиночества велика. Но главный источник этой неопределённости~--- допущение о возможном крайне маловероятном (трудношаговом) абиогенезе; наш отказ от трудных шагов обрезает этот хвост, неопределённость смещается к «жизнь не редка», и необходимость объяснять тишину возвращается~--- мы отвечаем на неё синхронностью. Наблюдательно модели различимы формой распределения времён выхода~--- узкий синхронный пик против широкого хвоста «трудных шагов».

%%% ===================================================================
\section{Критический путь из модели простых шагов}
\label{sec:steps}
%%% ===================================================================

\noindent\textit{Данный раздел даёт количественный вывод: формализует дисперсию времён выхода как концентрацию суммы большого числа независимых стадий и обосновывает универсальность параметров критического пути.}

\medskip
\noindent\textbf{Обозначения.} Везде ниже $\mu$~--- \emph{среднее} время элементарного шага по цепочке, $\mu \equiv \langle t_i\rangle = T/N$ (отдельные шаги могут иметь разные средние $\mu_i$; \S~\ref{sec:heterogeneity}).

\subsection{Допущения и концентрация суммы стадий}
\label{sec:concentration}

\noindent Модель опирается на явные предпосылки: (i)~независимость (некоррелированность) времён шагов; (ii)~конечность их дисперсий; (iii)~межпланетная универсальность \emph{статистики} шагов~--- одинаковые на всех планетах число шагов $N$ и распределение их времён (в частности, среднее время шага $\mu$),~--- обоснованная единственностью пути (\S~\ref{sec:uniqueness}) и плато скорости адаптации (\S~\ref{sec:resource_threshold}); шаги при этом \emph{не} предполагаются одинаковыми по длительности (\S~\ref{sec:heterogeneity}); (iv)~для \emph{числовой} инверсии в $\mu$~--- экспоненциальность одношагового ожидания ($c_v \approx 1$, \S~\ref{sec:inversion}), которая для самой концентрации не нужна. Приближённая нормальность распределения используется лишь в одном месте~--- для хвоста видимости в аргументе тишины (\S~\ref{sec:silence_sigma}).

\medskip
Постулат: эволюционное усложнение идёт $N$ малыми высоковероятными шагами (по Дарвину~\cite{Darwin1859}, \emph{natura non facit saltum}; не редкими «трудными»; \S~\ref{sec:planetary_steps}). Время выхода $T_{\mathrm{exit}} = \sum_{i=1}^{N} t_i$, шаги независимы. \textbf{Шаги не обязаны быть одинаковыми по длительности:} $i$-й шаг имеет, вообще говоря, своё среднее $\mu_i$ и свою дисперсию $\sigma_i^2$. Символ $\mu$ ниже всюду обозначает \emph{среднее} время шага $\mu \equiv \langle t_i\rangle = T/N$ (а не общую для всех шагов длительность), а $\sigma_{\mathrm{s}}^2 \equiv \langle\sigma_i^2\rangle$~--- среднюю пошаговую дисперсию. Ядро механизма~--- не форма распределения и не равенство шагов, а \textbf{аддитивность среднего и дисперсии} (тождество Бьенеме): для независимых шагов $\mathbb{E}[T_{\mathrm{exit}}] = \sum_i\mu_i = \mu N$ и $\mathrm{Var}(T_{\mathrm{exit}}) = \sum_i\sigma_i^2 = N\sigma_{\mathrm{s}}^2$, откуда $\sigma = \sigma_{\mathrm{s}}\sqrt{N}$, а \textbf{относительный разброс} $\sigma/T \propto 1/\sqrt{N}$ убывает с числом шагов. (Влияние \emph{неоднородности} $\mu_i$ и $\sigma_i$ на $\sigma$ разобрано отдельно в \S~\ref{sec:heterogeneity}: оно сводится к множителю порядка единицы, если ни один шаг не доминирует.) Для этого достаточно \emph{независимости} (некоррелированности) шагов и конечности их дисперсий~--- нормальность \emph{не} требуется. Вводя коэффициент вариации одного шага $c_v = \sigma_{\mathrm{s}}/\mu$, из $\sigma = \sigma_{\mathrm{s}}\sqrt{N}$ и $\sigma_{\mathrm{s}} = c_v\mu$ получаем общее соотношение $\sigma = c_v\mu\sqrt{N} = c_v\sqrt{\mu T}$, то есть $\sigma^2 = c_v^2\mu T$ и $\sigma/T = c_v/\sqrt{N}$. Частный случай $c_v = 1$ ($\sigma_{\mathrm{s}} \approx \mu$) отвечает экспоненциальной (пуассоновской) природе ожидания шага и даёт $\sigma = \mu\sqrt{N} = \sqrt{\mu T}$. Это \emph{отдельное} (параметрическое) допущение сверх аддитивности: на нём держится числовая инверсия в $\mu$, но \emph{не} сам вывод о концентрации, ведь $\sigma/T = c_v/\sqrt{N}$ убывает как $1/\sqrt{N}$ при любом конечном $c_v$. Значение $c_v \approx 1$ при этом не постулируется, а \emph{восстанавливается} из двух независимых каналов (\S~\ref{sec:mu_from_sigma}). Дополнительно при большом $N$ форма суммы приближается к нормальной (центральная предельная теорема, ЦПТ)~--- это понадобится лишь для оценки хвоста видимости в \S~\ref{sec:silence_sigma}, но не для самой концентрации. Подчеркнём: мы \emph{не} постулируем, что все шаги длятся одно время. Вместо прежнего постулата равенства шагов модель опирается на два робастных требования: (1)~велико \emph{число} шагов $N$; (2)~ни один отдельный шаг не доминирует в сумме $\sum_i\sigma_i^2$~--- нет неразложимого «трудного» шага с собственным временем порядка целевой $\sigma$ (\S~\ref{sec:heterogeneity}, \S~\ref{sec:planetary_steps}). Иными словами, важны \emph{число} шагов и \emph{время самого медленного} из них, а не равенство их длительностей. Среднее время шага $\mu$ при этом берётся сопоставимым для абиотической и биотической фаз~--- это упрощение: данные (часы Шарова, индекс сборки) указывают на разные характерные скорости усложнения (\S~\ref{sec:cosmic_ensemble}, \S~\ref{sec:open}), и его учёт сдвигает $\sigma$ лишь в пределах фактора $\sim1{,}5$ (\S~\ref{sec:heterogeneity}). В опорной форме аддитивность даёт $\sigma = \sqrt{\mu T}$ с $T \approx 13{,}8$~Gyr~--- полным временем от теплового гейта CMB (\S~\ref{sec:common_start}). В упрощённой модели настоящей работы межпланетная дисперсия определяется этим единственным механизмом; численные $\sigma$ и $\mu$ получаются из тишины (\S~\ref{sec:inversion}). Формула применяется напрямую, без поправок на ускорение, а само ускорение прогресса связано не с сокращением такта, а с ростом ступеньки и параллельной экспансией (\S~\ref{sec:cascade}); постоянство \emph{среднего} такта $\mu$ между фазами здесь~--- идеализация (см. оговорку выше и \S~\ref{sec:heterogeneity}), а не строго доказанный факт.

Малая межпланетная дисперсия $\sigma = \sqrt{\mu T}$ возникает из двух независимых источников. Первый~--- \textbf{общий старт}: тепловой гейт CMB задаёт единое начало критического пути для всей Галактики ($\sigma_{\mathrm{CMB}} \approx 0$, \S~\ref{sec:common_start}). Второй~--- \textbf{большое число малых шагов}: при $N \sim 10^{14}$ независимых шагов аддитивность дисперсии сжимает относительный разброс до $\sigma/T = c_v/\sqrt{N} \sim 10^{-7}$ (при $c_v \approx 1$). Из этих двух факторов следует, что момент выхода цивилизации определяется \emph{не} временем формирования планеты, а общим космологическим стартом и кинематикой цепочки шагов. Межпланетное постоянство числа шагов $N$ и \emph{среднего} такта $\mu$ (а не равенство отдельных шагов между собой)~--- следствие универсальности законов химии и физики, обеспечивающих статистически одинаковую цепочку повсюду.

\subsection{Чувствительность дисперсии к неоднородности шагов}
\label{sec:heterogeneity}

\noindent При независимости шагов дисперсия времени выхода равна сумме пошаговых дисперсий:
\begin{equation}
\sigma^2 = \mathrm{Var}(T_{\mathrm{exit}}) = \sum_{i=1}^{N}\sigma_i^2,
\label{eq:varsum}
\end{equation}
и это тождество \emph{не} использует равенства шагов. Сама $\sigma$ есть мера отклонения шагов от среднего: при детерминированных шагах ($\sigma_i=0$) было бы $\sigma=0$, то есть $\sigma\approx10^3$~лет~--- это и \emph{есть} следствие неоднородности, а не величина, которой неоднородность угрожает. Вопрос поэтому в чувствительности $\sigma$ к \emph{форме} распределения $\{\mu_i,\sigma_i\}$.

\noindent\textbf{Замкнутая формула.} Пусть шаг $i$ имеет среднее $\mu_i$ и разброс $\sigma_i = c_v\,\mu_i$ (общий коэффициент вариации внутри шага), а сами средние $\mu_i$ распределены по шагам со средним $\mu \equiv T/N$ и коэффициентом вариации $CV_\mu$ (мера того, \emph{насколько шаги различаются по длительности}). Тогда из~\eqref{eq:varsum} точно:
\begin{equation}
\boxed{\;\sigma = c_v\,\sqrt{1+CV_\mu^{2}}\;\sqrt{\mu\,T}\;}
\label{eq:sigma_general}
\end{equation}
Здесь два разных «отклонения от среднего». \emph{Разброс внутри шага} $c_v$ входит \textbf{линейно} (его значение $c_v\approx1$ восстанавливается из согласования каналов, \S~\ref{sec:mu_from_sigma}). \emph{Разброс средних между шагами} $CV_\mu$~--- ровно та неоднородность, что отдельные $\mu_i$ различаются вокруг среднего $\mu$,~--- входит под корнем как $\sqrt{1+CV_\mu^2}$, то есть \textbf{очень слабо}: при $CV_\mu=0$ имеем опорную форму $\sigma=\sqrt{\mu T}$; при $CV_\mu=1$ (разброс длительностей порядка самой длительности)~--- множитель лишь $\sqrt2\approx1{,}4$; даже $CV_\mu=10$ (стандартное отклонение в 10 раз больше среднего) даёт фактор $\sim10$. Таким образом, к замене «все шаги равны» на «$\mu$~--- среднее по цепочке» дисперсия \emph{устойчива}: умеренная неоднородность остаётся в пределах уже заявленной порядковой неопределённости (фактор $\sim1{,}5$).

\noindent\textbf{Единственное исключение~--- доминирующий шаг.} Слабость зависимости от $CV_\mu$ исчезает, когда сумма~\eqref{eq:varsum} начинает определяться \emph{одним} членом. Если один шаг имеет $\mu_{\max}\gg\mu$, то
\begin{equation}
\sigma \approx \sqrt{\big(\sqrt{\mu T}\,\big)^2 + (c_v\mu_{\max})^2} = \sqrt{(10^3~\text{лет})^2 + (c_v\mu_{\max})^2}:
\end{equation}
«море» из $N\sim10^{14}$ быстрых шагов даёт фиксированный пол $\sim10^3$~лет, а любой отдельный шаг добавляется в квадратурах. Порог, при котором один шаг сравнивается со всем остальным,
\begin{equation}
c_v\,\mu_{\max}\;\gtrsim\;\sigma\approx10^3~\text{лет}\quad\Longleftrightarrow\quad \mu_{\max}\gtrsim\sqrt{N}\,\mu\approx10^3~\text{лет}\approx10^{7}\mu .
\end{equation}
То есть отдельный шаг может быть в $\sim10^{7}$ раз медленнее среднего ($40$~мин $\to$ до $\sim10^3$~лет) и всё ещё не портить синхронность; но шаг с собственным временем заметно дольше $\sim10^3$~лет мгновенно задаёт $\sigma$ в одиночку. Эукариогенез как единый \emph{неразложимый} шаг ($\mu_{\max}\sim2$~Gyr) дал бы $\sigma\sim2\times10^9$~лет~--- концентрация уничтожена.

\noindent\textbf{Вывод.} Дисперсия $\sigma\approx10^3$~лет робастна к любой неоднородности \emph{быстрых} шагов и определяется по существу двумя величинами~--- числом шагов $N$ и временем самого медленного шага,~--- а не постулатом равенства длительностей. Концентрация рушится в единственном случае: если существует неразложимый «трудный» шаг с собственным временем $\gtrsim10^3$~лет, не ускоряемый параллелизмом популяции. Это переносит всю нагрузку на разложимость планетных шагов (\S~\ref{sec:planetary_steps}) и делает эукариогенез решающим тестом (\S~\ref{sec:eukaryo}).

\subsection{Единственность самого быстрого пути ($\Delta N = 0$)}
\label{sec:uniqueness}

\textbf{Термодинамическая сторона.} Само существование \emph{избытка} ресурсов (\S~\ref{sec:resource_threshold}) означает, что набор доступных путей диссипации \emph{конечен}: будь он неограничен, всегда нашёлся бы маршрут, которому не хватило бы и этого избытка, и никакого избытка тогда не возникало бы. На конечном наборе путей принцип максимального производства энтропии (MEPP~\cite{Martyushev,MartyushevSeleznev2006,Kleidon2010,Vallino2010}) отбирает \emph{наиболее диссипативный} маршрут; а поскольку рост сложности и есть рост производства энтропии (\S~\ref{sec:cascade}), самый диссипативный маршрут совпадает с самым \emph{быстрым} по усложнению. Иначе говоря, критический (самый быстрый) путь~--- это одновременно и самый \emph{диссипативный} путь. Будучи максимумом по скорости производства энтропии, он \emph{один} (точные совпадения скоростей имеют нулевую меру). Так существование единственного самого быстрого маршрута~--- критического пути~--- получает физическое основание, а не только постулируется. Ниже~--- кинематическая сторона.

Здесь важно оговориться, что критический путь по определению~--- не путь, привязанный к конкретной химии, а \textbf{максимум по скорости среди всех возможных маршрутов} усложнения, допускаемых физикой и химией данной Вселенной. До выхода на громкую фазу критический путь может быть только один: если бы существовала химическая основа, реализующая более быстрый маршрут, мы уже наблюдали бы его результат~--- громкую цивилизацию, выросшую на этом более быстром пути. Поскольку громких цивилизаций не наблюдается (космическая тишина), углеродно-органическая химия, на которой развивается жизнь нашего типа, совпадает с критическим путём данной Вселенной.

Теоретически возможны альтернативные формы жизни на иной химической основе, но они не образуют \emph{собственные} критические пути~--- по построению это просто более медленные маршруты, не достигающие статуса критического. Они также могут в конечном итоге приводить к возникновению разума и цивилизаций, но за существенно более длительные масштабы времени. Сам критический путь задаётся не выбором химии, а фундаментальными константами и параметрами Вселенной (скорость света, гравитационная и тонкая структурная постоянные, соотношения масс частиц и т.~д.): более быстрый критический путь в принципе мог бы существовать только в иной Вселенной с другими физическими константами.

\textbf{Кинематическая уникальность фиксирует число шагов ($N = \mathrm{const}$).} Уточним смысл слова \emph{шаг}: это элементарный такт усложнения (масштаба одного цикла репликации), а не крупная веха. То, что выглядит единичным «ключевым переходом» (эукариогенез, многоклеточность и т. п.), на критическом пути есть длинная цепочка таких шагов. Поэтому шагов очень много~--- $N \sim 10^{14}$ (\S~\ref{sec:inversion}),~--- но их число фиксировано. Критический путь определяется не минимумом \emph{числа} шагов, а \textbf{минимумом полного времени}: это самый короткий по времени путь к тому же уровню развития среди всех возможных альтернативных маршрутов~--- от пребиотической химии до заданного порога сложности (мы берём выход на громкую фазу~--- Тип~II по Кардашёву; \S~\ref{sec:inversion}). В рамках фундаментальных законов физики и химии такой самый быстрый маршрут кинематически уникален и может быть только один. Любые альтернативные маршруты существенно медленнее: на них шагов больше, сами шаги могут отличаться, а суммарное время заметно дольше. Поскольку самый быстрый путь физически уникален~--- задан фундаментальными константами Вселенной, а не выбором химии,~--- число шагов $N$ на нём жёстко зафиксировано и не имеет межпланетного разброса ($\Delta N = 0$).

\subsection{Время и форма цивилизации}
\label{sec:time_form}

Кинематическая уникальность (\S~\ref{sec:uniqueness}) относится к \emph{полному времени} $T_{\mathrm{exit}}$: среди всех траекторий к порогу сложности существует единственный глобальный минимум времени. Но проекция траектории на пространство \emph{форм цивилизации}~--- морфология биологического носителя, функциональный план, техносфера~--- высокоразмерна, а синхронизация задаёт лишь скалярный критерий. Концентрация суммы шагов сжимает межпланетный разброс в узкую «трубку» $\sigma \sim 10^3$~лет (\S~\ref{sec:concentration}, \S~\ref{sec:inversion}); в пространстве всех допустимых маршрутов, чьё $T_{\mathrm{exit}}$ попадает в эту трубку, могут лежать \emph{несколько} достаточно близких по длительности, но различных по форме цивилизации траекторий~--- вырождение у минимума времени, которое $\sigma$ не снимает. Отсюда: уникальность по времени не выбирает форму цивилизации.

Тем не менее модель (вместе с конвергентной эволюцией, \S~\ref{sec:ltee}) предсказывает не произвольную форму цивилизации, а воспроизводимый \emph{функциональный} план её биологического носителя~--- манипуляторные конечности, крупный социальный мозг, орудийность,~--- то есть гуманоида \emph{функционального типа}, а не принципиально иной эволюционно-морфологический базис. Вопрос другого порядка~--- \emph{насколько узка} морфологическая полоса носителя внутри этой «трубки»: достаточно ли общего гуманоидного \emph{типа} (разумные существа нашего функционального класса, но морфологически свои), или критический путь сжимает траекторию сильнее~--- к узкой полосе, практически неотличимой от \emph{Homo sapiens} и людей «как мы»? Модель последнего \emph{не} утверждает; этот спектр «функциональный тип $\leftrightarrow$ земная морфология вплотную» остаётся открытым (\S~\ref{sec:open}) и на него не отвечает ни синхронизация, ни $\sigma$~--- лишь эмпирическая конвергентность (\S~\ref{sec:ltee}), подтверждающая повторяемость \emph{функций}, но не задающая ни нижней, ни верхней границы близости к земному исходу.

\subsection{Универсальный такт $\mu$ и плато скорости ($\Delta\mu = 0$)}
\label{sec:plateau}

\textbf{Локальный избыток ресурсов и независимость от среды.} Под \emph{ресурсом} здесь понимается вся совокупность локальных условий, необходимых для шага усложнения~--- температура, кислотность, доступность нужных веществ, поток свободной энергии и т.~д.,~--- а не только энергетический поток или скорость производства энтропии. В этих терминах неоптимальные условия суть нехватка какого-либо из ресурсов, а избыток ресурсов означает оптимальные, идеальные условия. Прохождение критического пути требует ничтожно малого количества ресурсов по сравнению с общим пулом планеты (\S~\ref{sec:resource_threshold}), поэтому глобальные планетарные характеристики (размер, средняя температура, общая освещённость) его не лимитируют: для движения фронта достаточно сколь угодно малой локальной ниши с избытком ресурсов.

Ключевой момент~--- \textbf{насыщение скорости}. При избытке ресурсов (идеальных условиях) скорость прохождения шага выходит на плато: она перестаёт зависеть от условий и определяется уже только кинематикой химии~--- максимально быстрым физико-химически допустимым маршрутом. Это не отдельный постулат, а следствие двух независимых пределов, обоснованных (\S~\ref{sec:resource_threshold}): биохимического пола времени репликации и популяционно-генетического потолка скорости бесполой адаптации (клональная интерференция), который насыщается с ростом размера популяции выше скромного порога. Поэтому в идеальном оазисе скорость каждого шага задаётся не средой, а химией (для шага $i$~--- свой предел $\mu_{\min,i}$), и для каждого шага одинакова на всех планетах; речь о планетонезависимости \emph{потактовой} скорости, а не о равенстве разных шагов между собой. Планетам не нужно быть идеально одинаковыми: достаточно, чтобы на каждой большая площадь поверхности и совокупный ресурсный пул гарантировали в \emph{каждый момент времени} хотя бы одну нишу с избытком ресурсов на этом плато~--- а также возможность относительно быстрой экспансии между оазисами внутри планеты. Ресурсы в одном оазисе не избыточны для \emph{всех} стадий усложнения: для одних шагов нужные условия есть здесь, для других~--- в другом месте. Поэтому критический путь требует механизма подтягивания ресурсов с других оазисов (колонизация, миграция, транспорт метаболитов и энергии). Эволюция идёт длинными цепочками последовательных шагов усложнения, каждая из этих цепочек протекает там, где условия ближе всего к плато, в самом быстром месте. Разброс глобальных условий сдвигает лишь долю и размещение таких оазисов, но не $\mu$ на них, поэтому межпланетного разброса $\mu$ нет ($\Delta \mu = 0$).

Отсюда видно, почему здесь \emph{не} возникает «общей моды», разрушающей концентрацию (\S~\ref{sec:open}): общая мода появилась бы, если бы межпланетный разброс задавался \emph{условиями планеты} (температурой, составом, климатом), общими для всех её шагов и потому коррелирующими их между собой. Но критический путь по построению задаётся не условиями планеты, а \emph{универсальными свойствами химических элементов}: он всегда идёт при избытке ресурсов, на плато $\mu = \mu_{\min}$. Поэтому общепланетного множителя замедления не существует, и остаётся лишь \emph{независимая} пошаговая стохастика химической кинетики~--- ровно та, для которой работает аддитивность дисперсии и $\sigma/T = c_v/\sqrt{N}$. Это исключение общей моды справедливо \emph{при условии}, что на каждой планете в каждый момент есть хотя бы один оазис на плато (ничтожный ресурсный порог и быстрая внутрипланетная экспансия, \S~\ref{sec:resource_threshold}); само это условие однородности и есть нагруженное допущение модели (\S~\ref{sec:open}).

Таким образом, критический путь представляет собой универсальную кинематическую траекторию, заданную фундаментальными константами Вселенной и реализуемую~--- как показывает наблюдаемая тишина~--- через углеродно-органическую химию. В отсутствие межпланетного разброса параметров дисперсия времени выхода определяется исключительно стохастикой шагов через аддитивность дисперсии:
\begin{equation}
\boxed{\sigma = \mu\sqrt{N} = \sqrt{\mu T},}
\end{equation}
а все внешние факторы гетерогенности физически устранены (здесь записана опорная форма с $c_v \approx 1$; общий случай~--- $\sigma = c_v\sqrt{\mu T}$, \S~\ref{sec:inversion}).

\subsection{Каскадная модель: направленность и рост диссипации}
\label{sec:cascade}

Усложнение~--- каскад переходов по энергетическим ступенькам со свободной энергией Гельмгольца $\Delta F = \Delta U - T\Delta S$. Каждая ступенька~--- метастабильное состояние с памятью (время удержания $\tau \sim \exp(\Delta U/kT)$); производство энтропии эффективно понижает барьеры, делая каскадный подъём проходимым.

Сама по себе ступенчатая структура симметрична. \textbf{Направленность} возникает из асимметрии: с ростом сложности система растёт $\Rightarrow$ больше степеней свободы $\Rightarrow$ плотность микросостояний увеличивается $\Rightarrow$ переход «вправо» энтропийно выгоднее (фазовый объём более сложных состояний больше). Асимметрия проявляется не в \emph{темпе} шага, а в его \textbf{высоте}: каждый шаг вправо сопровождается бóльшим приростом энтропии $\Delta S$ и бóльшей диссипацией свободной энергии. Вместе с репликацией (храповик, \S~\ref{sec:ratchet}) это даёт статистический дрейф к усложнению вдоль стрелы времени (телеономия; ср.~Lotka~\cite{Lotka1922}, MEPP~\cite{Martyushev}, классическое термодинамическое обоснование~\cite{SchneiderKay1994}, диссипативная адаптация в самосборке~\cite{England2015}).

Растёт при этом прежде всего \textbf{высота ступеньки и суммарное производство энтропии} $\dot{S}_{\mathrm{prod}}$, а не частота тактов. $\dot{S}_{\mathrm{prod}}$ ускоряется «вправо» по двум причинам: (i) \textbf{экспансия}~--- репликация наращивает число систем на переднем фронте экспоненциально; (ii) более сложные метастабильные структуры требуют большего притока свободной энергии на единицу. Количественно эта тенденция фиксируется независимо измеряемой \emph{плотностью потока свободной энергии} $\Phi_m$ (эрг$\cdot$с$^{-1}\cdot$г$^{-1}$), которая растёт на порядки вдоль той же лестницы~--- звезда $\sim1$, планета $\sim10^2$, растение $\sim10^3$, животное $\sim10^4$, мозг $\sim10^5$, технологическое общество $\sim10^6$~\cite{Chaisson2011}. Само же \emph{среднее} время элементарного такта $\mu$ (один цикл репликации, один акт наследуемой памяти) остаётся \textbf{приблизительно постоянным} и снизу ограничено биохимией копирования наследуемой памяти~--- поэтому $\mu \approx$ время генерации (\S~\ref{sec:inversion}) и применима опорная форма $\sigma = \sqrt{\mu T}$ ($\mu$~--- среднее время шага, а не общая для всех шагов длительность; \S~\ref{sec:heterogeneity}). Это и есть упрощающее допущение настоящей работы: \textbf{среднее время шага $\mu$ примерно одно и то же на всём протяжении эволюции жизни} (отдельные шаги при этом различаются по длительности, что несущественно для $\sigma$, пока ни один не доминирует~--- \S~\ref{sec:heterogeneity}). Кажущееся ускорение прогресса на поздних стадиях обеспечивается не сокращением такта отдельного шага, а сменой носителя памяти (гены $\to$ культура $\to$ техника) и массовым параллелизмом (той же экспансией).

\subsection{Биологический храповик}
\label{sec:ratchet}

В одиночной механической системе видны и прямые, и обратные переходы. В биологии обратное движение (вымирание, редуктивная эволюция) идёт постоянно, но на уровне отдельных линий, и эти линии исчезают, не оставляя наблюдателя.

Ключевой усилитель~--- \textbf{репликация с памятью}: удачная ступенька копируется, и число систем на ней растёт экспоненциально. Чтобы откатилась вся популяция, независимый откат должен случиться у огромного числа копий разом~--- вероятность пренебрежимо мала. Репликация превращает слабую термодинамическую асимметрию в практически необратимый \textbf{храповик}~--- передний фронт сложности назад не откатывается (ср.~биологический закон необратимости Долло~\cite{Dollo1893,Gould1970}). Термодинамически этот храповик можно проиллюстрировать концепцией Джереми Ингланда: деление бактерии на две части~--- процесс относительно простой, но спонтанное слияние двух образовавшихся половин обратно в одну функциональную бактерию статистически и термодинамически невероятно из-за диссипации тепла и огромной разницы в энтропии состояний~\cite{England2013}.

\subsection{Разложение трудных шагов на простые}
\label{sec:planetary_steps}

На планетарном участке остаётся следующий риск. Длительные планетные эпохи (например, $\sim 2$~млрд лет до эукариот) сами по себе \textbf{не аномальны}: это просто \emph{длинные цепочки} обычных тактов (большое число шагов $N$), дающие тот же закон $\sigma = \sqrt{\mu T}$. Аномальный вклад в межпланетный разброс возник бы лишь от \textbf{неразложимого} медленного планетного шага~--- единичного перехода с собственным временем $\gtrsim10^3$~лет (порог доминирования из \S~\ref{sec:heterogeneity}), который нельзя представить цепочкой высоковероятных подшагов и который не является общим для планет.

Это, в иной форме,~--- гипотеза «трудных шагов» Картера~\cite{Carter1983}: редкие, по-настоящему неразложимые переходы с большой дисперсией времени ожидания. Но за более чем 40 лет обсуждения ни один конкретный шаг не доказан неразложимым; сама гипотеза опирается на допущения, которые недавний пересмотр называет сомнительно обоснованными как антропно, так и эволюционно~--- эффект «убранной лестницы» (успешные линии вытесняют конкурентов, маскируя их независимое происхождение), вытеснение нишы изменением среды и принципиальная ненаблюдаемость вымерших независимых линий~\cite{Mills2025}. Текущая тенденция исследований прямо противоположна: каждый из кандидатов на «трудный шаг» при ближайшем рассмотрении распадается на цепочку простых~--- происхождение эукариот через архей группы Asgard и их культивируемого представителя \emph{Prometheoarchaeum syntrophicum}~\cite{Spang2015,Zaremba2017,Imachi2020}, синтрофные и водородные гипотезы симбиогенеза~\cite{Martin1998,LopezGarcia2020}, актиновый цитоскелет у Lokiarchaeum~\cite{Rodrigues2023} и новые филогеномные реконструкции 2025--2026~гг.~\cite{Tobiasson2026,Speijer2025,BravoArevalo2025}.

Таким образом, вся нагрузка ложится на разложимость планетных шагов: пока не предъявлен ни один доказанно неразложимый шаг, постулат модели имеет эмпирическую поддержку, но не доказательство. Тест эукариогенеза (\S~\ref{sec:eukaryo}) остаётся решающим именно потому, что это последний серьёзный кандидат на «трудный шаг», для которого общепринятой механистической модели пока нет.

\subsection{Современные шаги~--- технологические, а не генетические}

На нынешнем этапе критический путь движется не медленной биологической эволюцией человека, а цивилизационно-технологическим прогрессом. При этом интернет, компьютер, авиация, искусственный интеллект~--- это не отдельные шаги, а \emph{вехи}: каждая сама складывается из огромного числа элементарных шагов усложнения (научных работ, инженерных улучшений, актов обмена знанием) и не возникает за один атомарный акт. Атомарный шаг на этом уровне намного мельче такой вехи. Это согласуется с \S~\ref{sec:cascade}: с ростом сложности меняется носитель шага (гены $\to$ культура $\to$ технологии), а кажущееся ускорение прогресса обеспечивается массовым параллелизмом и экспансией (ростом числа агентов и каналов передачи знания), а не сокращением такта отдельного элементарного шага. Темп выхода на стадию «громкости» за $\sim 10^3$~лет задаётся при этом динамикой энергопотребления (\S~\ref{sec:inversion}), а не изменением $\mu$.

\subsection{Ресурсный порог критического пути}
\label{sec:resource_threshold}

\textbf{Плато скорости и эталонный фронт.} Постоянство такта $\mu$ обосновывается не только энергетикой, но и популяционной генетикой. В бесполой популяции одновременно конкурируют многие полезные мутации: пока одна линия фиксируется, сопоставимые по приспособленности линии отбрасываются~--- это \emph{клональная интерференция}. Скорость адаптации растёт с численностью $N_c$ и потоком полезных мутаций $U_b$ лишь \textbf{логарифмически} и при больших $N_c$ \textbf{насыщается} (модель Герриша--Ленски~\cite{GerrishLenski1998}, travelling-wave теория~\cite{DesaiFisher2007,ParkKrug2007}; экспериментально~--- «предел скорости», не зависящий от потока мутаций~\cite{deVisser1999,Lang2013}). Насыщение наступает при $N_c\,U_b \gtrsim 1$, т.е.\ при $N_c \gtrsim 10^{6}$ клеток ($U_b \sim 10^{-6}$); дальнейший рост $N_c$ скорость почти не меняет~--- приспособленность не успевает передаваться через последовательные выметания. Снизу её ограничивает \emph{биохимический пол}~--- минимальное время деления ($\mu_{\min}\approx$ время генерации, $\sim 20$~мин у \emph{E.~coli}).

Для оценки энергобюджета критического пути берём \textbf{эволюционный фронт максимального темпа}~--- популяцию, уже находящуюся на этом плато. Эталон~--- одна колба долгосрочного эксперимента Ленски (LTEE~\cite{Lenski2017}): 10~мл, $\sim 5 \times 10^{8}$ клеток при суточном пике численности ($\sim 5 \times 10^{7}$~клеток/мл). Это на $\sim 2$--$3$ порядка \emph{выше} порога $N_c \sim 10^{6}$; эффективный размер $N_e \approx 3 \times 10^{7}$~--- тоже глубоко на плато. Именно поэтому LTEE наблюдает максимальный кинетический темп~--- все 12 линий идут почти одной кривой приспособленности~\cite{Wiser2013,Lenski1994},~--- и служит правильным масштабом «переднего фронта». Его энергетика:
\begin{itemize}
\item мощность $\approx 4{,}5 \times 10^{-5}$~Вт; $\dot{S}_{\mathrm{front}} \approx 1{,}5 \times 10^{-7}$~Вт/К.
\end{itemize}
Планета (Земля): $\dot{S}_{\mathrm{planet}} \approx 4{,}6 \times 10^{14}$~Вт/К; $\dot{S}_{\mathrm{bio}} \approx 4 \times 10^{11}$~Вт/К; прокариот $\sim 5 \times 10^{30}$.

Отношение фронт/планета: $\dot{S}_{\mathrm{front}}/\dot{S}_{\mathrm{planet}} \approx 3 \times 10^{-22}$; клетки $\approx 10^{-22}$; площадь~--- фронту достаточно $\sim 2$~см$^2$ ($\sim 3 \times 10^{-19}$ поверхности планеты).

\textbf{Вывод.} Порог выхода на критический путь $\sim 10^{-19}$--$10^{-22}$ доли планеты~--- ничтожен. Это снимает \emph{энергетическое и ресурсное} лимитирование пути глобальными условиями планеты; кинетическое постоянство такта ($\mu = \text{const}$) следует из насыщения скорости на плато (выше).

\textbf{Следствие для $\mu\approx\mathrm{const}$.} Именно насыщающаяся, логарифмическая зависимость делает $\mu$ робастным к межпланетным вариациям размера популяции $N_c$: чтобы сдвинуть $\mu$ вдвое, $N_c$ должно измениться на много порядков. Порог $\sim 10^{6}$ клеток ничтожен~--- даже минимальный оазис $\sim 2$~см$^2$ содержит $\sim 10^{8}$ клеток, поэтому любая планета с жизнью уже на плато $\mu_{\min}$, и $N_c$-зависимость межпланетного разброса не создаёт ($\Delta\mu = 0$).

\subsection{Локальные оазисы против глобального средового гейтирования}

Концепция критического пути вступает в содержательную дискуссию с гипотезой «глобального средового гейтирования» (environmental gating), отстаиваемой, в частности, в работе Mills et al.~\cite{Mills2025}. Согласно их взглядам, крупные эволюционные переходы (эукариогенез, многоклеточность) задерживались не из-за внутренней сложности, а в ожидании созревания планеты (достижения глобального порога кислорода, климатической стабилизации и т.д.). Такой механизм неизбежно приводил бы к значительному разбросу времён выхода цивилизаций, поскольку планетные таймеры (скорость накопления $O_2$) у всех планет уникальны.

В рамках нашей модели геология и планетарные таймеры не являются сдерживающими шлюзами (гейтами), благодаря двум эмпирически подтвержденным факторам:
\begin{enumerate}
    \item \textbf{Кислородные оазисы.} Биохимическая эволюция не ждет глобального изменения планеты, а протекает в локальных очагах с избытком ресурсов. Геологические данные показывают, что мелководные «кислородные оазисы», создаваемые локальными сообществами цианобактерий, существовали сотни миллионов лет до Великого окислительного события (GOE) в условиях полностью аноксичной планеты. Эволюционные шаги, требующие кислорода, шли внутри этих локальных оазисов на максимальной кинетической скорости, а последующее глобальное насыщение планеты привело лишь к пространственному расселению уже сформированных форм жизни.
    \item \textbf{Тест порядка появления (Mills~2014 против Mills~2025).}  Гипотеза гейтирования утверждает, что появление животных жестко лимитировалось ростом кислорода. Однако независимые измерения потребности ранних многоклеточных (на примере современных губок) показывают, что они способны выживать при концентрациях кислорода всего $0{,}5$--$4{,}0\%$ от современного уровня~\cite{Mills2014}. Такие концентрации были достигнуты на Земле задолго до реального появления животных. Тот факт, что животные возникли со значительным временным лагом после прохождения этого порога, свидетельствует, что кислород не был лимитирующим стоп-фактором. Эволюция шла со своим собственным кинетическим темпом биологических переходов.
\end{enumerate}

\subsection{Границы применимости: локальные оазисы и планетный масштаб}
\label{sec:applicability_bounds}

Тезисы о локальном избытке ресурсов (\S~\ref{sec:plateau}, \S~\ref{sec:resource_threshold}) и механизме локальных оазисов (выше) полностью снимают проблему планетарного гейтирования для \emph{ранних} этапов критического пути~--- микробных, биохимических и одноклеточных: для них достаточно сколь угодно малой ниши с избытком ресурсов, и эволюция не ждёт глобального созревания планеты.

По мере движения «вправо» по каскаду усложнения (\S~\ref{sec:cascade}) объём потребляемых ресурсов и энергопотребление монотонно растут (растёт высота ступеньки и производство энтропии), а кажущийся темп эволюции ускоряется за счёт экспансии. Поэтому тезис «локального оазиса достаточно» справедлив в первую очередь для шагов \emph{до} уровня порядка Типа~0{,}5 по шкале Кардашёва. Уже на нашем уровне ($K \approx 0{,}73$) цивилизация требует существенной доли планетарного энергетического пула; дальше потребление будет только нарастать.

Для \emph{поздних} шагов, требующих глобальных условий, планетный фактор может оставаться значимым: например, возникновение крупных активных хищников с высоким метаболизмом, которые физически не могут существовать в границах локального оазиса~\cite{Knoll2014}. В таких редких звеньях критический путь переходит на планетный масштаб~--- это естественная граница применимости модели: $\mu \approx \mathrm{const}$ и независимость от глобальной среды хорошо обоснованы для подавляющего большинства шагов и времени пути, но не для всех поздних переходов без оговорок.

\subsection{Экспериментальное свидетельство воспроизводимости пути}
\label{sec:ltee}

Долгосрочный эксперимент Ленски (LTEE)~\cite{Lenski2017}~--- лучший доступный полигон: 12 параллельных «проигрываний плёнки» из одного предкового клона, более 73\,000 поколений. Это контролируемая постановка спора Конвей-Морриса~\cite{ConwayMorris2003} $\leftrightarrow$ Гулда~\cite{Gould1989}.

Показатель синхронности~--- не идентичность мутаций, а \textbf{повторяемость функционального результата и времени} между независимыми линиями. Все 12 линий идут почти одной кривой роста приспособленности~\cite{Wiser2013,Lenski1994}. На уровне результата эволюция повторяема~--- прямое свидетельство того, что самый быстрый критический путь усложнения фиксирован и воспроизводим.

Контингентные сингулярности существуют, но вне критического пути: аэробный рост на цитрате (Cit$^+$) возник лишь в 1 из 12 линий~\cite{Blount2008,Blount2012}. Это медленный детур к побочному ресурсу; критический путь обходит такие ответвления. Cit$^+$ показывает, что контингентные переходы существуют, но не лежат на критическом пути, пока есть более быстрый маршрут.

\textbf{Оговорка о масштабе.} LTEE меряет микроэволюцию; перенос микро$\to$макро~--- отдельное допущение.

%%% ===================================================================
\section{Эволюция до Земли: космическая фаза критического пути}
\label{sec:cosmic_ensemble}
%%% ===================================================================

\noindent\textit{Данный раздел прослеживает критический путь до Земли: общий старт от реликтового гейта, «часы сложности», локально конвергентную пребиотическую фазу, преклеточный порог, выше которого усложнение продолжается лишь на планете, и переход на планетарную когорту.}

\subsection{Общий старт и реликтовый гейт}
\label{sec:common_start}

Критический путь начинается не с появления жизни на планете и не с образования Земли, а с первого момента, когда жидкостная химия стала возможна во Вселенной. Этот момент задаётся \textbf{тепловым гейтом реликтового излучения}: при $100 \lesssim (1+z) \lesssim 137$ температура CMB составляла 273--373~K, делая жидкостную химию возможной на любых каменистых телах независимо от наличия звезды, уже через $\sim 10$--17~млн лет после Большого взрыва~\cite{Loeb2014}. Реликт изотропен до $\Delta T/T \sim 10^{-5}$, поэтому тепловой гейт открывается синхронно по всей Вселенной с разбросом $\Delta t \sim 10^{-5} \cdot t \sim 10^2$--$10^3$~лет~--- того же порядка, что $\sigma \approx 10^3$~лет, независимо полученная из тишины (\S~\ref{sec:inversion}). Синхронный старт критического пути обусловлен, таким образом, не биологией, а космологией: часы запущены одновременно для всей Вселенной. В масштабах Галактики ($\sim 30$~кпк, или $\sim 0{,}3$~кпк физического размера при $z\sim 120$) флуктуации CMB подавлены шелковским затуханием до уровней $\ll 10^{-5}$, поэтому $\sigma_{\mathrm{CMB}}$ в галактическом масштабе $\approx 0$: реликтовый гейт обеспечивает \emph{общий старт} для всех систем Галактики, не внося межсистемного разброса. Тепловой гейт обеспечивает \emph{общий старт} для всей Галактики; единственный вклад в межпланетную дисперсию даёт стохастика суммы шагов (аддитивность дисперсии), применённая ко всей цепочке усложнения~--- от этого старта до технологической цивилизации.

\subsection{Часы сложности}
\label{sec:complexity_clocks}

Наличие «часов сложности» следует из критического пути: при кинематически уникальной траектории (\S~\ref{sec:steps}) и каскадном закреплении в наследуемой памяти (\S~\ref{sec:cascade}) накопленная сложность растёт монотонно и примерно экспоненциально, поэтому любая монотонная её мера~--- хронометр после калибровки темпа; «часы» здесь не дополнительный постулат.

\textbf{Первые часы: функциональный геном (Шаров).} Алексей Шаров предложил использовать рост функционального (неизбыточного) размера генома как «часы» происхождения и эволюции жизни~\cite{Sharov}: регрессия по времени даёт экспоненциальный рост $\sim$ в $6{,}3$ раза за 1~млрд лет (удвоение $\sim$ каждые 376~млн лет), а экстраполяция к одному нуклеотиду переносит зарождение жизни на $\sim 9{,}7 \pm 2{,}5$~млрд лет назад~\cite{Sharov2013}~--- ещё до образования Земли. Отсюда следует, что критический путь начинается не на планете: значительная часть усложнения~--- от пребиотической химии до первого репликатора~--- протекает в открытом космосе, на миллиарды лет раньше появления океанов. Ключевая идея настоящего раздела: пребиотическая химия протекает \textbf{конвергентно} вблизи каждой звёздной системы~--- одни и те же законы физики и химии, одни и те же исходные элементы дают одинаковый результат повсюду, не требуя обмена веществом между разными частями галактики.

\textbf{Независимые «часы сложности»: индекс сборки.} Теория сборки (Assembly Theory) предлагает независимую меру молекулярной сложности~--- \emph{индекс сборки} (MA), равный минимальному числу операций копирования, необходимых для построения молекулы из атомов~\cite{Marshall2021,Cronin2023}. Опубликованные значения: глицин MA$\approx$5, аденин MA$\approx$8, АТФ MA$\approx$13, типичные метаболиты MA$\approx$12--15. Шкала $\log_2(\mathrm{MA})$ совместима со шкалой Шарова: для случайных полимерных последовательностей MA$\approx N$ (длина цепи), поэтому $\log_2(\mathrm{MA}) \approx \log_2 N$~--- именно то, что измеряет Шаров. Это открывает возможность \textbf{унификации} двух шкал в единые «часы сложности» от Большого взрыва до наших дней: абиотическая фаза (химический рост MA) плавно переходит в биологическую (геномный рост по Шарову). Количественная оценка показывает, что обе кривые совместимы, однако скорость усложнения в абиотической фазе ($\sim 0{,}6$~Gyr$^{-1}$ в логарифмической шкале) примерно в 4--5 раз ниже, чем в биологической ($\sim 2{,}7$~Gyr$^{-1}$ по Шарову). Настоящая работа \emph{не} ставит задачи построения полной унифицированной шкалы и использует упрощённую однофазную модель суммы стадий; разработка расширенных часов сложности вынесена в открытые вопросы (\S~\ref{sec:unified_clocks}).

\textbf{Третьи «часы»: плотность потока свободной энергии (Чейссон).} Сложность можно отслеживать не только информационно (геном, индекс сборки), но и \emph{энергетически}. Эрик Чейссон предложил \emph{плотность потока свободной энергии} $\Phi_m$ (эрг$\cdot$с$^{-1}\cdot$г$^{-1}$) как универсальную метрику сложности и движущий фактор эволюции~\cite{Chaisson2011}: она монотонно растёт вдоль всей космической лестницы~--- звёзды $\sim 1$, планеты $\sim 10^2$, растения $\sim 10^3$, животные $\sim 10^4$, мозг $\sim 10^5$, технологическое общество $\sim 10^6$~--- и, отложенная против времени появления структур, даёт примерно экспоненциальный рост на протяжении $\sim 13$~млрд лет (удвоение порядка $\sim 0{,}5$--$1$~млрд лет). В отличие от двух предыдущих мер, эти часы фиксируют не информацию, а \emph{диссипацию}, и потому напрямую сопрягаются с каскадной термодинамической картиной (\S~\ref{sec:cascade}): каждый шаг «вправо» поднимает высоту ступеньки и $\Phi_m$. \emph{Уточнение:} из механики критического пути (\S~\ref{sec:cascade}) физически растёт прежде всего \textbf{удельная скорость производства энтропии} $\dot S_{\mathrm{prod}}$, а не энергия как таковая,~--- именно $\dot S_{\mathrm{prod}}$ напрямую связана со ступеньками сложности (каждый шаг «вправо» увеличивает прирост энтропии $\Delta S$). Поток свободной энергии $\Phi_m$~--- удобная и измеримая, но по сути \emph{производная} величина; при фиксированной $T$ они пропорциональны ($\dot S = \Phi/T$), так что для роли часов это одно и то же, но физически обоснована именно скорость производства энтропии. Именно эти часы использованы ниже для энергетической оценки уровня C* из бюджета пылинки (\S~\ref{sec:idp_dust}, \S~\ref{sec:freeze}).

\textbf{К объединению трёх часов.} Три независимые меры~--- функциональный геном (Шаров), индекс сборки MA (Cronin/Walker) и плотность потока свободной энергии $\Phi_m$ (Чейссон)~--- описывают одну траекторию усложнения с разных сторон: информационной (первые две) и энергетической (третья). Их объединение в единые «часы сложности» от Большого взрыва до цивилизации выглядит естественным: все три растут примерно экспоненциально и связаны каскадной моделью (\S~\ref{sec:cascade}), где рост наследуемой информации ($N$, MA) и рост диссипации ($\Phi_m$)~--- две стороны одного подъёма по энергетическим ступенькам с закреплением результата в памяти. Согласование темпов нетривиально: информационные часы в биофазе идут в $\sim 4$--$5$ раз быстрее, чем в абиотической, а $\Phi_m$ растёт даже при почти неизменном геноме (у прокариот геном вышел на плато, тогда как диссипация продолжала расти за счёт метаболической оптимизации). Но это не противоречие: все три часов меряют \emph{одно и то же}~--- рост сложности вдоль критического пути,~--- лишь с немного разных сторон (наследуемая информация против диссипации), и потому \textbf{дополняют друг друга}. Их расхождения не обесценивают часы, а несут дополнительный смысл~--- указывают на смену ведущего носителя усложнения; вместе три меры дают более полную картину единой траектории, чем любая по отдельности. Построение согласованных трёхкомпонентных часов (с переводом между $\log N$, $\log \mathrm{MA}$ и $\log \Phi_m$ и явными пофазовыми темпами) мы оставляем отдельной задачей (\S~\ref{sec:unified_clocks}); для настоящей работы существенно лишь то, что все три независимо указывают на единый, примерно экспоненциальный и стартующий ещё до образования Земли критический путь.

\subsection{Локальная конвергентность пребиотической химии}
\label{sec:local_conv}

\textbf{Порог C*}~--- уровень преклеточной сложности, ниже которого открытая космическая эволюция устойчива, а выше которого необходимы устойчивая жидкофазная среда, ресурсный пул, концентрированный поток свободной энергии, активный трансмембранный транспорт и полноценный метаболизм~--- возможные лишь на тёплой планете. C* соответствует РНК-подобным репликаторам или протоклеткам без полноценной мембраны, но \emph{не} полноценной клетке уровня JCVI-syn3.0 (473~гена; \S~\ref{sec:freeze}).

Синхронность достижения порога C* вблизи разных звёздных систем обеспечивается не обменом веществом между частями галактики~--- такой обмен работает на временных масштабах $\gg \sigma$ и физически не мог бы объяснить синхронизацию,~--- а \textbf{локальной конвергентностью} пребиотической химии. Вблизи каждой формирующейся звёздной системы пребиотические реакции протекают на одних и тех же физических входных данных:

\begin{itemize}
    \item \textbf{Химический состав:} межзвёздная среда галактического диска содержит примерно одинаковое соотношение органогенных элементов (H, C, N, O, S, P), задаваемое нуклеосинтезом предшествующих поколений звёзд. Разброс по диску для данной эпохи~--- не более $\sim 2$--$3$ раз.
    \item \textbf{Источники энергии:} ультрафиолетовое излучение, космические лучи, ударные волны~--- универсальны и присутствуют вблизи каждой формирующейся системы.
    \item \textbf{Физические условия:} температуры молекулярных облаков и протопланетных дисков определяются одними и теми же физическими законами и схожи повсюду.
\end{itemize}

Одни и те же исходные элементы плюс одни и те же законы физики и химии порождают \textbf{конвергентный} пребиотический путь: как и в биологической эволюции (\S~\ref{sec:ltee}), самый быстрый химический путь к заданному уровню сложности воспроизводим. Каждая звёздная система \emph{локально} и \emph{независимо} «открывает» одни и те же органические молекулы в одной и той же последовательности~--- аминокислоты, нуклеооснования, рибозу, простые репликаторы~\cite{Furukawa2019,Callahan2011,Oberg2021}. Разброс в скорости определяется химической кинетикой, а не случайными особенностями конкретной системы.

\subsection{Заморозка космической фазы и порог клеточной сложности C*}
\label{sec:freeze}

Конвергентная пребиотическая химия (\S~\ref{sec:local_conv}) работает до тех пор, пока шаги требуют мало энергии и не предъявляют жёстких требований к среде. Что именно останавливает дальнейшее усложнение выше порога C* в открытом космосе, \emph{не установлено}; ниже~--- \emph{рабочая гипотеза} модели. Она не привязана к конкретному носителю~--- справедлива для любой открытой космической площадки (пыль, кометы, метеориты и т.~д.), а не только для мелкой пыли (\S~\ref{sec:idp_dust}). Мы предполагаем, что барьер связан прежде всего с \emph{условиями среды}, а не с глобальным «остыванием» Вселенной и дефицитом свободной энергии в межзвёздной среде: космическое остывание идёт слишком медленно, а энергию пребиотическая химия на таких носителях черпает в основном из локальных источников~--- УФ молодых звёзд и ударных процессов (\S~\ref{sec:local_conv}),~--- а не из реликтового фона, задающего лишь ранний тепловой гейт (\S~\ref{sec:common_start}). По этой гипотезе поддержание липидной мембраны, активный трансмембранный транспорт и полноценный метаболизм требуют устойчивой жидкофазной среды (при типичной химии~--- жидкой воды при нужном давлении и температуре), ресурсного пула и длительного концентрированного потока свободной энергии. Обеспечить это на любом открытом носителе в вакууме трудно. В её рамках эволюция в открытом космосе «замораживается» на некотором \textbf{пороговом уровне сложности} C*, соответствующем преклеточным структурам: именно поэтому в метеоритах и протопланетных дисках до сих пор обнаруживают лишь пребиотическую органику~--- аминокислоты, нуклеооснования, рибозу~--- но не живые клетки. Шаги выше C* переносятся на тёплые планеты с атмосферой, гидросферой и ресурсным пулом, необходимым для дальнейшей эволюции. Количественную оценку положения C* из бюджета одной пылинки см.~\S~\ref{sec:cstar_dust}.

\textbf{Оценка момента достижения C*.} Минимально возможный геном самовоспроизводящейся клетки определён экспериментально: синтетическая клетка JCVI-syn3.0~\cite{Hutchison2016} содержит 473~гена при длине генома $\sim 531$~т.п.н.~--- нижняя эмпирическая граница клеточной сложности. По кривой Шарова этот порог был достигнут
\begin{equation}
\Delta t_{\text{C*}} = \tau \cdot \log_2\!\left(\frac{G_{\text{C*}}}{G_0}\right),
\end{equation}
где $\tau = 376$~Мyr~--- время удвоения функционального генома~\cite{Sharov2013}, $G_0 \sim 5$--$40$~т.п.н.~--- эффективная сложность первого репликатора (не экстраполяционная точка «1 нуклеотид», а порядок величины минимального самовоспроизводящегося полимера), $G_{\text{C*}} \approx 531$~т.п.н.~--- сложность минимальной клетки (JCVI-syn3.0). Подстановка даёт $\Delta t_{\text{C*}} \approx 5{,}1$--$6{,}3$~Gyr от начала репликации; калибровка по ранним бактериям Земли ($\sim 3{,}5$~млрд лет назад) согласуется с тем же диапазоном, что отвечает дате достижения C*~--- примерно $3{,}4$--$4{,}6$~млрд лет назад. Этот диапазон \textbf{совпадает с возрастом Земли} (4,5~млрд лет): к моменту образования Земли преклеточная космическая эволюция как раз «дозрела» до уровня, позволяющего быстрый переход к полноценной клетке на планете. Это объясняет, почему жизнь на Земле появилась столь быстро после образования океанов~--- уже готовая преклеточная химия на уровне C* была доставлена на готовых носителях~--- пыли, кометах и метеоритах протопланетного диска (\S~\ref{sec:local_conv}),~--- и до первой клетки оставался минимальный шаг.

\subsection{Предполагаемый носитель~--- мелкая пыль}
\label{sec:idp_dust}

Где именно аккумулируется преклеточная сложность~--- на поверхности планеты, в межзвёздной и межпланетной пыли (IDPs), на кометах или метеоритах~--- строго не установлено (\S~\ref{sec:open}). Среди этих носителей мы склоняемся к \textbf{мелкой пыли} как основной площадке космической фазы критического пути~--- по двум причинам. \emph{Во-первых}, пылинок чрезвычайно много, и, в отличие от крупных метеоритов, мелкие частицы ($\sim 10$~мкм) тормозятся в атмосфере плавно и оседают на поверхность планеты \emph{без существенного нагрева} (ниже температуры пиролиза органики $\sim 600^\circ$C), доставляя накопленную сложную химию неразрушенной; именно такая пыль вносит основную долю внеземного органического углерода на Землю~\cite{Anders1989,ChybaSagan1992,Flynn2004}. \emph{Во-вторых}, энергетический бюджет критического пути не только ничтожен, но и \emph{убывает к более ранним стадиям}: плотность потока свободной энергии растёт вдоль лестницы усложнения (\S~\ref{sec:cascade}), поэтому пребиотическая химия требует энергии \emph{ещё меньше}, чем микробный фронт ($\sim 10^{-5}$~Вт, \S~\ref{sec:resource_threshold}), и столь малый бюджет физически реализуем даже на отдельной пылинке~--- ей не нужны ни планетарный ресурсный пул, ни близкая звезда. Оба довода согласуются с тем, что в метеоритах и протопланетных дисках находят именно \emph{преклеточную} органику (аминокислоты, нуклеооснования, рибозу), но не клетки (\S~\ref{sec:freeze}).

\subsection{Оценка C* из бюджета одной пылинки}
\label{sec:cstar_dust}

Потребуем, чтобы критический путь укладывался в энергетику \emph{одной} пылинки~--- то есть чтобы отдельное зерно служило самодостаточным оазисом (\S~\ref{sec:idp_dust}). Для зерна радиусом $r \sim 5$~мкм ($\sim 10$~мкм в диаметре) на $\sim 1$~а.е. от формирующейся звезды при потоке $F \sim 10^3$~Вт/м$^2$ получаем перехваченную мощность $P \approx F\pi r^2 \sim 10^{-7}$~Вт; доля фотохимически активного УФ $f_{\mathrm{UV}} \sim 10^{-2}$ даёт $P_{\mathrm{UV}} \sim f_{\mathrm{UV}} P \sim 10^{-9}$~Вт; производство энтропии $\dot S \approx P/T \sim 10^{-10}$~Вт/К при $T \sim 300$~K. Сравним это с энергобюджетом фронта максимального темпа адаптации из \S~\ref{sec:resource_threshold} (колба LTEE, $\sim 45$~мкВт), внеся две поправки. \emph{Во-первых}, сами $45$~мкВт~--- с запасом: колба ($5\times10^{8}$~клеток) лежит на $\sim 2{,}5$ порядка выше порога выхода на плато скорости ($N_c\sim 10^{6}$, \S~\ref{sec:resource_threshold}), так что минимальный фронт, ещё держащийся на плато, требует лишь $\sim 10^{-7}$~Вт. \emph{Во-вторых}, на преклеточной стадии удельная диссипация на $\sim 2$ порядка ниже современной микробной (рост $\Phi_m$ по Чейссону за $\sim 4$~млрд лет, \S~\ref{sec:cascade}), что снижает фронт ещё на два порядка~--- до $\sim 10^{-9}$~Вт.

Итог: минимальный фронт критического пути на уровне C* требует $\sim 10^{-9}$~Вт~--- ровно порядка полезного УФ-бюджета одной пылинки. Условие «достаточно одной пылинки» тем самым выполняется и \emph{независимо фиксирует} положение C*: оно лежит примерно на $4$--$5$ порядков ниже фронта LTEE ($\approx 2{,}5$ из них~--- запас по численности, $\approx 2$~--- более ранняя, менее диссипативная химия), то есть заведомо \emph{ниже} минимальной клетки~--- на преклеточном уровне репликаторов/протоклеток. Это~--- независимое (энергетическое) подтверждение гипотезы (\S~\ref{sec:freeze}), что космическая фаза замерзает именно на преклеточном C*, а доставка готового C* на пылинках объясняет быстрый старт жизни на Земле.

\subsection{Когорта земляподобных планет на критическом пути}
\label{sec:bimodal}

Наличие допланетарного участка критического пути (до порога C*, \S~\ref{sec:freeze}) определяет, какие именно планеты входят в текущий синхронизированный пик.

\textbf{Когорта критического пути} объединяет земляподобные планеты, сформировавшиеся \emph{до} достижения C* ($\sim 3{,}4$--$4{,}6$~млрд лет назад). Эти планеты получили преклеточную химию уровня C* непосредственно с образованием и без разрыва запустили планетарный участок. Стартовые условия для всех них являются common-mode: химия уровня C* в межзвёздной среде достигнута локально и независимо (\S~\ref{sec:local_conv}) вблизи каждой системы. Межпланетный разброс определяется стохастикой суммы шагов: $\sigma \approx \sqrt{\mu\,T} \sim 10^3$~лет. Это синхронизированный пик, к которому принадлежим мы.

Земля (возраст 4,5~млрд лет) находится вблизи \textbf{верхней возрастной границы когорты}: она образовалась почти в момент достижения C*, то есть является одной из самых «молодых» планет в этом пике. Планеты, сформировавшиеся \emph{после} C*, имеют временной разрыв в непрерывности критического пути: в период между образованием планеты и текущим моментом они не находились на критическом пути эволюции.

Для оценки $n$: в текущий синхронизированный пик входят лишь земляподобные планеты старше $\sim 4$~млрд лет (п.~4.3), что дополнительно ограничивает $n \sim 10^4$--$10^5$ (\S~\ref{sec:inversion}).

%%% ===================================================================
\section{Оценка \texorpdfstring{$\sigma$}{σ} из тишины}
\label{sec:inversion}
%%% ===================================================================

\noindent\textit{Данный раздел инвертирует космическую тишину в оценку межпланетной дисперсии $\sigma$ времён выхода на детектируемую («громкую») стадию и числа $n$ планет Галактики, на которых критический путь идёт прямо сейчас; сопоставляет результаты с независимой астрофизикой и микробиологией и формулирует следствия для SETI.}

\subsection{Время до громкой стадии}
\label{sec:loud_time}

В этой традиции~--- и в согласии с разделением на «громкие» и «тихие» цивилизации в недавних моделях~\cite{Hanson2021}~--- мы отождествляем «громкость» с выходом на Тип~II по шкале Кардашёва ($P_{\mathrm{II}} = 10^{26}$~Вт по классической интерполяции Сагана~\cite{Kardashev1964,Sagan1973}): энергобюджет уровня звезды, дающий детектируемую на межзвёздных и межгалактических расстояниях сигнатуру (тепловое излучение сфер Дайсона и других мегаструктур~\cite{Dyson1960,Wright2014}). Это не новое предположение настоящей работы, а принятый в SETI-литературе порог детектируемости.

Текущая мощность цивилизации $P_0 \approx 2 \times 10^{13}$~Вт ($\approx 20$~ТВт на 2024 год~\cite{EI2024}), что по шкале Сагана соответствует $K \approx 0{,}73$~\cite{Sagan1973}. При экспоненциальном росте $3\%$ в год время до Типа~II:
\begin{equation}
t_{\mathrm{loud}} = \frac{\ln(P_{\mathrm{II}}/P_0)}{\ln(1{,}03)} \approx 990~\text{лет}.
\end{equation}
Исторический средний темп роста мирового энергопотребления за последние десятилетия составляет $\approx 1{,}5\%$--$2{,}0\%$ в год (по данным IEA~\cite{IEA2023} и EI~\cite{EI2024}), поэтому значение $3\%$ используется здесь как реалистично-оптимистичный сценарий для ускоряющейся стадии развития техносферы. Чувствительность этой величины к темпу роста представлена в табл.~1.

\begin{center}
\textbf{Таблица 1.} Время до перехода к стадии громкости ($t_{\mathrm{loud}}$) при различной скорости экспоненциального роста ($P_{\mathrm{II}} = 10^{26}$~Вт, $P_0 = 2 \times 10^{13}$~Вт)

\medskip
\begin{tabular}{c c c}
\hline
Темп роста в год & до Типа~I ($10^{16}$~Вт) & до Типа~II ($= t_{\mathrm{loud}}$) \\
\hline
1{,}5\% & 417 лет & 1964 лет \\
2{,}0\% & 314 лет & 1477 лет \\
2{,}5\% & 252 лет & 1184 лет \\
\textbf{3{,}0\%} & \textbf{210 лет} & \textbf{990 лет} \\
3{,}5\% & 181 лет & 850 лет \\
4{,}0\% & 158 лет & 746 лет \\
\hline
\end{tabular}
\end{center}

\subsection{Аргумент тишины как ограничение на дисперсию}
\label{sec:silence_sigma}

Эмпирический факт: детектируемых техносигнатур не наблюдается. Систематические поиски и статистические аргументы парадокса Ферми показывают, что отсутствие сигналов накладывает верхнюю границу на долю цивилизаций, достигших уровня, при котором они становились бы межзвёздно заметными~\cite{Annis1999,Wright2014,Hanson2021}. Если бы межпланетный разброс времён выхода $\sigma$ существенно превышал $t_{\mathrm{loud}}$, заметная доля близких цивилизаций опережала бы нас более чем на $t_{\mathrm{loud}}$, успела бы стать громкой и была бы видна.

Количественно примем реалистичную модель~--- \emph{толстый галактический диск}: $n$ миров равномерно в диске радиуса $R = 5\times10^4$~св.~лет и полной толщины $h = 10^3$~св.~лет (объём $V = \pi R^2 h$). В \textbf{базовой} постановке (табл.~2) наблюдатель в центре диска, $c = 1$~св.~год/год; поправка при $R_\odot$ и неоднородном профиле звёзд~--- после табл.~2. Мир на трёхмерном расстоянии $s = \sqrt{r^2+z^2}$ ($r$~--- радиус в плоскости диска от центра, $z$~--- высота) виден сегодня, если опередил нас не меньше чем на $t_{\mathrm{loud}} + s$~--- постройку громкой фазы ($t_{\mathrm{loud}}$) плюс ход света ($s/c$),~--- с вероятностью $\Phi(-(t_{\mathrm{loud}}+s)/\sigma)$ (верхний хвост нормального распределения времён выхода~--- следствие ЦПТ). Ожидаемое число видимых соседей:
\begin{equation}
\boxed{\;\mathbb{E}[V] = \frac{n}{V}\int_{-h/2}^{h/2}\!\!\int_{0}^{R} \Phi\!\left(-\frac{t_{\mathrm{loud}}+\sqrt{r^2+z^2}}{\sigma}\right) 2\pi r\,\mathrm{d}r\,\mathrm{d}z \;\lesssim\; 1.\;}
\end{equation}
Конечная толщина существенна: на малых расстояниях $s \lesssim h/2$ соседи распределены трёхмерно, на больших~--- планарно (тонкий диск~--- частный предел $h\to 0$). Из $\mathbb{E}[V]\lesssim 1$ для каждого $n$ находят \textbf{наибольшее} $\sigma$, при котором тишина ещё возможна (верхняя граница, табл.~2); $n$~--- число планет Галактики, на которых критический путь идёт прямо сейчас, а $\sigma$~--- разброс именно до стадии детектируемости, а не до появления жизни.

\noindent\textit{Оговорка о нормальности.} Именно здесь~--- и только здесь~--- модель использует \emph{форму} распределения (гауссов $\Phi$), а не одну лишь концентрацию. Это самостоятельное допущение, причём задействуется оно в \emph{верхнем хвосте}, где сходимость суммы к нормали по центральной предельной теореме наиболее медленная: реальные крупные уклонения могут отклоняться от гауссовых. Поэтому оценка видимости (и численные границы табл.~2) модельно-зависима сильнее, чем сам вывод о концентрации $\sigma/T = c_v/\sqrt{N}$, который держится только на аддитивности дисперсии.

\noindent\textit{Примечание.} Интеграл берётся в замкнутом виде. С $a=t_{\mathrm{loud}}/\sigma$ и $\beta=h/2\sigma$ переход к 3D-расстоянию $s$ (при $R\to\infty$, поправка $\sim e^{-(R/\sigma)^2/2}$ пренебрежима) даёт $\mathbb{E}[V]=(n/V)\,4\pi\sigma^3\,W(a,\beta)$, $W(a,\beta)=\int_0^\beta K(a;y)\,dy$, где $K(a;y)=\int_y^\infty x\Phi(-(a+x))dx=\tfrac12[(1+a^2-y^2)\Phi(-(a+y))+(y-a)\varphi(a+y)]$. Внешний интеграл~--- конечная комбинация $\Phi$ и $\varphi$ в точках $a$ и $a+\beta$ (гауссовы моменты). Предельные случаи: $\beta\to0$ даёт 2D-форму $\sigma^2 K(a)$, $K(a)=\tfrac12[(1+a^2)\Phi(-a)-a\varphi(a)]$; $\beta\to\infty$~--- 3D-форму $\sigma^3 M(a)$, $M(a)=\tfrac13[(a^2+2)\varphi(a)-a(3+a^2)\Phi(-a)]$.

\begin{center}
\textbf{Таблица 2.} Верхняя граница $\sigma$ при $\mathbb{E}[V]=1$ (толстый диск: $R = 5\times10^4$, $h = 10^3$, $t_{\mathrm{loud}} = 990$~св.~лет). $n$ миров равномерно в объёме, наблюдатель в центре; $\sigma$~--- решение замкнутой формы из примечания. $\ell = \Gamma(4/3)\,(3 R^2 h/(4n))^{1/3}$~--- среднее 3D-расстояние до ближайшего соседа ($\rho = n/V$, $V = \pi R^2 h$; множитель $\pi$ в числителе и знаменателе сокращается).

\medskip
\small
\begin{tabular}{c c c}
\hline
$n$ & $\sigma$, лет & $\ell$, св. лет \\
\hline
$10^2$ & 7840 & 2370 \\
$10^3$ & 2980 & 1100 \\
$10^4$ & 1390 & 510 \\
\textbf{$3{,}8 \times 10^4$} & \textbf{990} & \textbf{330} \\
$10^5$ & 810 & 240 \\
$10^6$ & 570 & 110 \\
$10^7$ & 440 & 50 \\
$10^8$ & 360 & 25 \\
\hline
\end{tabular}
\normalsize
\end{center}

\noindent\textit{Поправка модели} (Monte Carlo, $N_{\mathrm{MC}} = 4\times10^5$): перенос наблюдателя на $R_\odot = 8$~кпк и замена однородности на $\propto e^{-r/R_d}$, $R_d = 3$~кпк \emph{ужесточают} ограничение тишины~--- верхняя граница $\sigma$ смещается вниз на $\sim 1$--$2\%$ по всему диапазону $n$ (не более $\sim 3\%$); для обратной задачи при $\sigma = 990$~лет допустимое $n$ ниже примерно на $\sim 5$--$6\%$ ($\approx 3{,}6\times10^4$ вместо $\approx 3{,}8\times10^4$).
\subsection{Число планет на критическом пути}
\label{sec:kepler}

В распределение входят не цивилизации как таковые, а «чашки Петри»~--- планеты, на которых может идти быстрый критический путь, дозревающий до сложной жизни и цивилизации. Их число даёт свежий астрофизический расчёт.

\textbf{Оценка Scherf--Lammer 2024.} Работа~\cite{Scherf2024,Lammer2024} оценивает максимальное число \emph{земных местообитаний} (Earth-like Habitats)~--- скалистых планет в зоне обитаемости сложной жизни, способных удержать устойчивую азот-кислородную атмосферу земного типа с малым $\mathrm{CO_2}$,~--- с учётом истории звездообразования, начальной функции масс, галактической зоны обитаемости и термической устойчивости атмосфер. Результат: во всём Млечном Пути таких сред \textbf{не более $n_{\text{EH}} \sim 6\times10^4$--$2{,}5\times10^5$}, преимущественно у G- и K-звёзд. Эта оценка уже включает ключевые срезы~--- возраст планеты, поколение и металличность звезды~\cite{Lineweaver2001,Behroozi2015}, положение в галактической зоне обитаемости и стабильность атмосферы. Полный набор критериев отбора: скорость звездообразования и начальная функция масс, галактическая зона обитаемости, преимущественно G- и K-звёзды, зона обитаемости \emph{сложной} жизни, частота скалистых планет ($\eta_\oplus$), одновременное наличие океанов и суши, вероятное требование крупной луны и термическая устойчивость N$_2$--O$_2$-атмосферы при заданной доле CO$_2$.

\textbf{Совпадение с тишиной.} Независимый результат из космической тишины (\S~\ref{sec:silence_sigma}, табл.~2)~--- $n \sim 10^4$--$10^5$ цивилизационных планет при $\sigma \sim 0{,}7$--$1{,}3$~тыс.~лет~--- совпадает с астрофизической оценкой Scherf--Lammer ($n_{\text{EH}} \sim 10^5$). Два независимых канала~--- астрофизика населённости и статистика молчания~--- сходятся на $n \sim 10^5$. Будь планет на критическом пути на порядки больше, тишина была бы нарушена (\S~\ref{sec:silence_sigma}). Поскольку $n_{\text{EH}}$~--- верхняя оценка числа пригодных сред, а не гарантированных цивилизаций, фактическое $n$ на критическом пути может быть ниже; это только ужесточает согласие с тишиной.

\subsection{Оценка длительности элементарного шага}
\label{sec:mu_from_sigma}

При постулате об универсальности и постоянстве параметров критического пути ($\sigma = \sqrt{\mu T}$) наблюдательная граница $\sigma$ (из условия тишины) \textbf{инвертируется} в оценку длительности элементарного шага:
\begin{equation}
\boxed{\mu = \frac{\sigma^2}{T}.}
\end{equation}
При $T = 1{,}38 \times 10^{10}$~лет для каждой строки табл.~2 (тишина при $\mathbb{E}[V]=1$) получаем согласованные параметры; при отождествлении $\sigma \approx t_{\mathrm{loud}}$ темп экспоненциального роста энергопотребления восстанавливается из формулы табл.~1. Число эволюционных шагов~--- $N = (T/\sigma)^2$ (табл.~3).

\noindent\textit{Логика вывода условна и однонаправлена.} Постулаты (модель шагов, $\mu \approx \mathrm{const}$) плюс модель $t_{\mathrm{loud}}$ дают границу $\sigma$ из тишины, а та~--- оценку $\mu = \sigma^2/T$. Число шагов $N = T/\mu$ при этом \emph{не} измеряется независимо, а следует из $\mu$ (это не замкнутый круг, но и не независимое подтверждение). Поэтому совпадение $\mu$ с временем деления клетки~--- \emph{фальсифицируемое предсказание двух каналов} (астрономия $\leftrightarrow$ микробиология), а не доказательство.

\noindent\textit{Чувствительность к форме распределения шага.} Без гипотезы $c_v = 1$ общее соотношение $\sigma^2 = c_v^2\mu T$ (\S~\ref{sec:cascade}) даёт $\mu = \sigma^2/(c_v^2 T) = (40~\text{мин})/c_v^2$: при $c_v \ne 1$ оценка $\mu$ смещается как $1/c_v^2$ (например, $c_v = 0{,}5 \Rightarrow \mu \approx 150$~мин; $c_v = 2 \Rightarrow \mu \approx 9$~мин), а $N = c_v^2 (T/\sigma)^2$. Однако $c_v$ не обязан задаваться извне: соотношение $\sigma^2 = c_v^2\mu T$ связывает три величины, и любые две задают третью. Подставив \emph{независимо} известные $\mu \approx 40$~мин (время деления, микробиология) и $\sigma \approx 990$~лет (тишина, астрономия), получаем $c_v^2 = \sigma^2/(\mu T) \approx 0{,}93$, то есть $c_v \approx 0{,}97 \approx 1$. Тем самым экспоненциальность одношагового ожидания не вводится произвольно, а \emph{следует} из согласования двух каналов: совпадение из «двухканального» становится трёхканальным (астрономия $\leftrightarrow$ микробиология $\leftrightarrow$ форма распределения шага). Сама концентрация $\sigma/T = c_v/\sqrt{N}$ от значения $c_v$ не зависит.

\begin{center}
\textbf{Таблица 3.} Сводная оценка: тишина (табл.~2), $t_{\mathrm{loud}}$ (табл.~1) и параметры критического пути ($T = 1{,}38 \times 10^{10}$~лет, $\mu = \sigma^2/T$). $n$~--- число планет на критическом пути в Галактике; $N$~--- число элементарных шагов; темп роста~--- скорость экспоненциального роста $P(t)$, при которой $t_{\mathrm{loud}} = \sigma$.

\medskip
\small
\begin{tabular}{c c c c c c}
\hline
$n$ & $\sigma$, лет & $\ell$, св.~лет & рост, \%/год & $\mu$, мин & $N$ \\
\hline
$10^2$ & 7840 & 2370 & 0{,}37 & 2340 & $3{,}1 \times 10^{12}$ \\
$10^3$ & 2980 & 1100 & 0{,}99 & 338 & $2{,}1 \times 10^{13}$ \\
$10^4$ & 1390 & 510 & 2{,}1 & 74 & $9{,}9 \times 10^{13}$ \\
\textbf{$3{,}8 \times 10^4$} & \textbf{990} & \textbf{330} & \textbf{3{,}0} & \textbf{37} & \textbf{$1{,}9 \times 10^{14}$} \\
$10^5$ & 810 & 240 & 3{,}7 & 25 & $2{,}9 \times 10^{14}$ \\
$10^6$ & 570 & 110 & 5{,}3 & 12 & $5{,}9 \times 10^{14}$ \\
$10^7$ & 440 & 50 & 6{,}9 & 7{,}4 & $9{,}8 \times 10^{14}$ \\
$10^8$ & 360 & 25 & 8{,}5 & 4{,}9 & $1{,}5 \times 10^{15}$ \\
\hline
\end{tabular}
\normalsize
\end{center}

\noindent На шкале Земли критический путь провёл в одноклеточной форме $\sim 3$~млрд лет (от $\approx 3{,}7$ до $0{,}6$~млрд лет назад~\cite{Nutman2016,Knoll2014}). При постулате $\mu \approx \mathrm{const}$ доля шагов равна доле времени, и из $N = T/\mu \approx 1{,}9\times10^{14}$ (табл.~3) следует: космическая дофаза пребиотического усложнения ($\sim 9$~Gyr) $\sim 65\%$, ранняя планетная и переходная фаза $\sim 9\%$, земная одноклеточная фаза $\sim 3{,}1/13{,}8 \approx 22\%$, многоклеточная и цивилизационная вместе $\sim 4\%$ (сумма $\approx 100\%$). Вклад \emph{поздних} носителей памяти (многоклеточность, культура, техника) в число шагов невелик ($\sim 4\%$). Значение $\mu \approx 40$~мин фиксируется именно микробной фазой; собственный такт пребиотической дофазы может отличаться от $40$~мин~--- это открытый вопрос (\S~\ref{sec:open}), и при его учёте $\sigma$ и инвертированное $\mu$ сдвигаются в пределах уже заявленной порядковой неопределённости (фактор $\sim 1{,}5$).

\noindent Полученная инверсией оценка $\mu \approx 40$~мин (в опорной точке табл.~3~--- $\approx 37$~мин; величина порядковая, интервал правдоподобия $25$--$75$~мин~--- \S~\ref{sec:intervals}) лежит в диапазоне генерационного времени микроорганизмов (\emph{E.~coli} $\sim 20$--$60$~мин). Космическая тишина (астрономия) и время деления клетки (микробиология) указывают на одну и ту же фундаментальную «единицу такта» эволюции. При ином темпе роста (столбец «рост» в табл.~3) значение $\mu$ вышло бы из биологически осмысленного окна. \textbf{Оговорка о выборе строки.} Подчеркнём, что из всех строк табл.~3 биологически осмысленна лишь узкая~--- $\mu \approx 20$--$60$~мин, отвечающая $n \sim 10^4$--$10^5$ и темпу роста энергопотребления $\sim 2$--$3\%$/год. Опорную строку ($n \approx 3{,}8\times10^4$, рост $3\%$) мы выбираем не как единственно возможную, а как \emph{наиболее правдоподобную}: темп $3\%$~--- реалистично-оптимистичный сценарий для ускоряющейся техносферы, а полученное $\mu \approx 37$~мин совпадает с измеренным временем деления микроорганизмов. Совпадение поэтому не выводит $\mu$ из первых принципов, а \emph{отбирает} из таблицы ту строку, где три независимых ориентира~--- микробиология ($\mu$), астрофизическая оценка населённости (Scherf--Lammer, $n$) и реалистичный темп роста энергопотребления~--- указывают на одно и то же значение. Это усиливает оценку, но не доказывает её. Через $\sigma = \sqrt{\mu T}$ то же окно $\mu \approx 20$--$60$~мин отвечает $\sigma \approx 730$--$1260$~лет; по табл.~2 такой диапазон $\sigma$ допускает при наблюдаемой тишине $n \sim 10^4$--$10^5$ цивилизационных планет в Млечном Пути. Это совместное фальсифицируемое предсказание двух независимых каналов: если населённость Галактики окажется $n \gg 10^5$, то либо $\sigma < 730$~лет (тогда $\mu$ быстрее времени деления клетки~--- физически сомнительно), либо тишина обязана нарушиться.

\subsection{Интервалы неопределённости и устойчивость $\sigma$}
\label{sec:intervals}

\noindent\textbf{Характер интервала.} Приводимый ниже интервал~--- \emph{не} частотный доверительный интервал (здесь нет выборочной статистики), а \emph{диапазон правдоподобия}, полученный распространением неопределённости входных параметров через соотношения $\mu=\sigma^2/(c_v^2T)$ и $N=c_v^2(T/\sigma)^2$. Он \emph{условен}: задаётся при населённости $n\in[10^4,10^5]$ (пересечение астрофизической оценки Scherf--Lammer, \S~\ref{sec:kepler}, и допустимого тишиной диапазона, табл.~2) и $c_v\approx1$ (\S~\ref{sec:mu_from_sigma}).

\noindent\textbf{Слабая зависимость $\sigma$ от $n$~--- ключевое свойство.} Граница $\sigma(n)$ из тишины (табл.~2) \emph{сильно сублинейна}: при изменении населённости в \textbf{10 раз} (от $10^4$ до $10^5$) $\sigma$ меняется лишь в $\approx1{,}7$ раза (с $1390$ до $810$~лет), тогда как $\mu$ и $N$~--- в $\approx2{,}9$ раза (как $\sigma^2$). Поэтому даже грубая неопределённость населённости \emph{не} расшатывает оценку $\sigma$: дисперсия времён выхода зафиксирована на уровне $\sim10^3$~лет почти независимо от того, сколько именно планет идёт по критическому пути. Это же делает робастным предсказание $\mu\sim40$~мин~--- оно остаётся в окне деления микроорганизмов ($20$--$60$~мин) при любой правдоподобной $n$.

\begin{center}
\textbf{Таблица 4.} Интервалы правдоподобия при $n\in[10^4,10^5]$, $c_v\approx1$, $T=1{,}38\times10^{10}$~лет.

\medskip
\begin{tabular}{l c c c}
\hline
Параметр & Центр & Диапазон ($n\in[10^4,10^5]$) & Главный источник \\
\hline
$\sigma$ & $1000$~лет & $800$--$1400$~лет & населённость $n$ \\
$\mu$ & $40$~мин & $25$--$75$~мин & $\sigma$ (как $\sigma^2$) \\
$N$ & $1{,}9\times10^{14}$ & $(1{,}0$--$2{,}9)\times10^{14}$ & $\sigma$ (как $\sigma^{-2}$) \\
$\ell$ & $330$~св.~лет & $240$--$510$~св.~лет & населённость $n$ \\
\hline
\end{tabular}
\end{center}

\noindent\textbf{Бюджет неопределённостей.} Сверх параметрического интервала (фактор $\sim2$, от $n$) действуют три систематики, которые мы оговариваем отдельно, чтобы не смешивать разнородные источники:
\begin{itemize}
\item \emph{Многофазность часов сложности} (\S~\ref{sec:open}): абиотическая и биотическая фазы имеют разные характерные скорости; учёт сдвигает $\sigma$ и $\mu$ в пределах фактора $\sim1{,}5$.
\item \emph{Коэффициент вариации шага} $c_v$: зафиксирован $\approx1$ согласованием каналов (\S~\ref{sec:mu_from_sigma}); при его отпускании $\mu\propto1/c_v^2$ (например, $c_v=0{,}7\Rightarrow\mu$ вдвое больше). В интервал не включён во избежание двойного счёта.
\item \emph{Гауссов хвост в аргументе тишины} (\S~\ref{sec:silence_sigma})~--- \textbf{доминирующая} и принципиально \emph{односторонняя} неопределённость: негауссовы крупные уклонения способны сдвинуть допустимое $n$ на порядки. Через слабую $\sigma(n)$ сама $\sigma$ при этом смещается скромнее (фактор $\sim2$--$3$), но $N$ и вывод о населённости для SETI~--- сильнее. Эта систематика не выражается симметричным $\pm$ и в табл.~4 не включена.
\end{itemize}
Расширение населённости до верхней оценки Scherf--Lammer ($n\approx2{,}5\times10^5$) сдвигает интервал к $\sigma\approx700$~лет, $\mu\approx19$~мин (нижний край микробного окна), $N\approx3{,}8\times10^{14}$, $\ell\approx175$~св.~лет~--- то есть остаётся в тех же порядках.

\subsection{Следствие для SETI}
\label{sec:seti}

Статистическая синхронизация меняет постановку вопроса. Обычно молчание Ферми трактуют в одной из двух крайностей: либо мы одни (или крайне редки), либо другие цивилизации значительно старше и обогнали нас. Модель предлагает третий ответ: \textbf{большинство цивилизаций возникают примерно одновременно}, поскольку концентрация суммы большого числа шагов сжимает разброс до $\sigma \sim 10^3$~лет~--- пренебрежимо мало на масштабе жизни Вселенной, так что ни одна цивилизация не оказывается значимо «старше» прочих. Отсюда другая интерпретация SETI: другие технологические цивилизации, если они есть, с высокой вероятностью находятся на \emph{сопоставимой} стадии~--- то есть являются \textbf{потенциальными конкурентами}, а не маяками или покровителями; молчание объясняется не отсутствием или удалённостью других, а тем, что все они появились недавно (по космическим меркам) и ещё не вышли на детектируемый уровень. Стратегии SETI, ориентированные на поиск сигналов значительно более развитых цивилизаций, могут оказаться принципиально неверными.

\begin{equation}
\boxed{\text{Тишина} \;\Longrightarrow\; \sigma \sim 800\text{--}1400~\text{лет} \;\Longrightarrow\; \mu \approx 25\text{--}75~\text{мин} \;\Longrightarrow\; n \sim 10^4\text{--}10^5,\;\ell \sim 10^2\text{--}10^3~\text{св.~лет}.}
\end{equation}

\noindent Здесь $n$~--- число планет Галактики, на которых критический путь идёт сейчас (не гарантированные цивилизации), $\ell$~--- среднее 3D-расстояние до ближайшей такой планеты (табл.~2--3); при согласованной точке $n \approx 3{,}8 \times 10^4$, $\sigma \approx t_{\mathrm{loud}}$ получаем $\ell \approx 330$~св.~лет~--- в согласии с оценкой Scherf--Lammer $n_{\text{EH}} \lesssim 10^5$ (\S~\ref{sec:kepler}).

%%% ===================================================================
\section{Предсказания и фальсифицируемость}
\label{sec:eukaryo}
%%% ===================================================================

Модель проверяема сразу на нескольких независимых фронтах; ниже~--- единый список предсказаний, каждое с явным условием фальсификации.
\begin{itemize}
\item \textbf{Конвергентность (палеобиология).} Крупные переходы к сложности конвергентны, повторяемы, фиксированы по числу и порядку. \emph{Фальсификация:} богатая контингентность по Гулду~--- если бы «переигрывание плёнки» (ср. LTEE, \S~\ref{sec:ltee}) систематически давало радикально иные исходы.

\item \textbf{Эукариогенез (решающий тест).} Эукариотическая клетка возникла на Земле, по-видимому, однократно за $\sim 2$~млрд лет~--- главный кандидат на уникальный медленный «трудный шаг» и потому самый острый тест модели. Контртезис: это не скачок, а длинная синтрофная цепочка высоковероятных подшагов~--- асгардские археи с эукариотическими сигнатурами~\cite{Spang2015,Zaremba2017}, культивируемый \emph{Ca.}~Prometheoarchaeum syntrophicum (модель E3: Entangle--Engulf--Endogenize)~\cite{Imachi2020}, синтрофная и водородная гипотезы~\cite{LopezGarcia2020,Martin1998}, актиновый цитоскелет Lokiarchaeum~\cite{Rodrigues2023}, доминантный вклад Asgard-архей и пошаговый переход FECA$\to$LECA~\cite{Tobiasson2026,Speijer2025,BravoArevalo2025}. Тогда длительность $\sim 2$~млрд лет~--- длина цепочки, а не один трудный шаг; полной механистической модели пока нет, поэтому тест остаётся \emph{открытым}. \emph{Фальсификация:} если хотя бы одно звено окажется устойчиво неразложимым, уникальным и медленным~--- не проходящим быстрее $\sim 10^3$~лет даже с учётом всей популяции параллельно эволюционирующих линий (параллельный поиск его не ускоряет),~--- это даёт большой разброс времён выхода и разрушает синхронность.

\item \textbf{Конвергентная форма цивилизации.} Критический путь задаёт самую быструю траекторию через пространство сложности \emph{по полному времени} (\S~\ref{sec:uniqueness}); вместе с конвергентной эволюцией (\S~\ref{sec:ltee}) это влечёт предсказание не только синхронности времён, но и сходимости \emph{формы цивилизации} (\S~\ref{sec:time_form}): технологические цивилизации на других планетах сходятся к воспроизводимому \emph{функциональному} плану биологического носителя~--- манипуляторные конечности, крупный социальный мозг, орудийность,~--- то есть к гуманоидам нашего типа, а не к рептильно-динозавровым или принципиально иным сценариям.

\emph{Подтверждение:} обнаружение на другой планете технологической цивилизации с носителем функционально гуманоидного типа~--- и при этом \emph{никого} существенно более развитого~--- подтвердило бы конвергенцию формы цивилизации и синхронность времён ($\sigma \sim 10^3$~лет, когда ни одна цивилизация не ушла вперёд на космологические сроки). Именно совмещение этих двух условий~--- «та же функциональная форма цивилизации» и «никто не опередил»~--- отличает наше предсказание от сценариев, где иные цивилизации либо принципиально иной формы, либо значительно старше и развитее нас (\S~\ref{sec:seti}). \emph{Фальсификация:} обнаружение технологической цивилизации на принципиально ином эволюционно-морфологическом базисе; обнаружение цивилизации, существенно более развитой, чем наша, разрушило бы тезис синхронности.

\item \textbf{Астрономическое (синхронность и тишина).} Высокая синхронность появления цивилизаций ($\sigma \sim 10^3$~лет); при наблюдаемой тишине это ограничивает населённость~--- $n \sim 10^4$--$10^5$ планет (ближайший сосед $\sim 10^2$--$10^3$~св.~лет). Инверсия $\sigma$ даёт средний такт $\mu = \sigma^2/T$ (\S~\ref{sec:inversion}); согласование с биологическим масштабом~--- отдельная проверка двух каналов, а не следствие единого шага. \emph{Фальсификация:} разброс $\sigma \sim 10^8$--$10^9$~лет (космологическое одиночество) либо населённость $n \gg 10^5$ при сохраняющейся тишине.

\item \textbf{Возраст планеты и стадия жизни.} Достигнутая стадия монотонно растёт со временем, проведённым на критическом пути, поэтому коррелирует с возрастом планеты. На планетах \emph{старше} Земли (когорта, стартовавшая без разрыва, \S~\ref{sec:bimodal}) жизнь, если она есть, уже не может находиться на простой одноклеточной стадии: за время не меньшее земного она прошла её и достигла сложных, технологических форм. В силу синхронности (\S~\ref{sec:seti}) такие планеты не на эпохи впереди, а в пределах $\sigma \sim 10^3$~лет от нашего уровня~--- но заведомо \emph{за} микробной фазой. Напротив, простая одноклеточная жизнь ожидается \emph{только} на планетах, существенно \emph{моложе} Земли~--- сформировавшихся уже после достижения порога C* и потому отставших от синхронизированного пика (\S~\ref{sec:bimodal}), которые сейчас находятся в ранней фазе пути. \emph{Фальсификация:} обнаружение богатой \emph{простой} (микробной, доцивилизационной) биосферы на планете заметно старше Земли~--- без признаков продвижения к сложности~--- противоречило бы монотонности и синхронности критического пути.

\item \textbf{Подъём по Кардашёву к Типу~II (наше будущее).} Выход на «громкую» стадию идёт через устойчивый экспоненциальный рост энергопотребления (исторически $\sim 1{,}5$--$3\%$ в год) и достигает Типа~II за $t_{\mathrm{loud}} \sim 10^3$~лет (\S~\ref{sec:loud_time}). Модель трактует это как \emph{типовой} исход критического пути и тем самым предсказывает наше собственное будущее~--- продолжение подъёма Тип~I\,$\to$\,II на горизонте $\sim 10^3$~лет (одновременно посылка о «громкости» цивилизаций и проверяемое во времени следствие). \emph{Фальсификация:} устойчивое отклонение траектории от такого роста~--- долгое плато, спад, «великий фильтр» ниже Типа~II у нас и у наблюдаемых цивилизаций~--- лишает силы $t_{\mathrm{loud}}$, инверсию «тишина $\to \sigma$» и саму посылку аргумента тишины.
\end{itemize}
Таким образом, предложенный подход фальсифицируем независимо по нескольким фронтам~--- палеобиологии, астрономии и динамике техносферы; ближайший решающий тест~--- разложимость эукариогенеза.

%%% ===================================================================
\section{Открытые вопросы}
\label{sec:open}
%%% ===================================================================

\begin{itemize}
\item \textbf{Независимость и однородность шагов (главные допущения).} Аддитивность дисперсии (подавление разброса как $1/\sqrt{N}$) требует \emph{независимости} времён шагов: при общей моде~--- когда «медленная» планета медленна на каждом шаге~--- в пределе полной корреляции $\sigma \propto N$, а не $\sqrt{N}$, и концентрация рушится. В нашей модели общая мода исключена самим построением критического пути: он задаётся универсальными свойствами химических элементов, а не условиями планеты, и всегда идёт при избытке ресурсов на плато $\mu = \mu_{\min}$, так что общепланетного множителя замедления нет (\S~\ref{sec:resource_threshold}). Однако это исключение справедливо лишь \emph{при условии} непрерывного наличия хотя бы одного оазиса на плато; именно это условие межпланетной однородности и остаётся нагруженным. Поэтому независимость шагов и межпланетная универсальность их \emph{статистики}~--- одинаковые на всех планетах число шагов $N$ и распределение времён шагов (но \emph{не} равенство отдельных шагов между собой, \S~\ref{sec:heterogeneity})~--- суть главные нагруженные допущения; они защищаются физически (\S~\ref{sec:resource_threshold}, \S~\ref{sec:planetary_steps}), но из первых принципов не доказаны. Отметим, что к \emph{внутрипутевой} неоднородности длительностей $\sigma$ устойчива (\S~\ref{sec:heterogeneity}): нагружено именно требование однородности \emph{между} планетами, а не равенство шагов внутри пути. Остаточный риск~--- неразложимый медленный планетный шаг с собственным временем $\gtrsim10^3$~лет (\S~\ref{sec:heterogeneity}, \S~\ref{sec:planetary_steps}; тест \S~\ref{sec:eukaryo}).

\item \textbf{Конвергентность формы цивилизации.} Спектр «гуманоид функционального типа $\leftrightarrow$ \emph{Homo sapiens} вплотную» и границы, которые на него накладывает «трубка» $\sigma$, разобраны в \S~\ref{sec:time_form}. Кратко: модель предсказывает функциональный план гуманоидного носителя цивилизации (\S~\ref{sec:time_form}), но не утверждает совпадения с земной морфологией; открытый вопрос~--- \emph{насколько близка} форма цивилизации по второй оси. Эмпирическая опора~--- лишь конвергентная эволюция (\S~\ref{sec:ltee}); фальсификация~--- \S~\ref{sec:eukaryo}, пункт о форме цивилизации.

\item \textbf{Нормальность в хвосте тишины.} Приближённая нормальность распределения времён выхода используется только в аргументе тишины (\S~\ref{sec:silence_sigma}), причём в хвосте, где сходимость к нормали наиболее медленная; это самостоятельное и более уязвимое допущение, чем сама концентрация.

\item \textbf{Порог громкости и $t_{\mathrm{loud}}$.} Порог «громкости» (Тип~II) заимствован из SETI-литературы~\cite{Annis1999,Wright2014,Hanson2021}; вывод $t_{\mathrm{loud}}$ зависит от модели энергопотребления (\S~\ref{sec:loud_time}).

\item \textbf{Старт критического пути.} Синхронный старт задаётся тепловым гейтом CMB и материальным гейтом первых звёзд~\cite{Naoz2006} (\S~\ref{sec:cosmic_ensemble}); оба ранние и почти универсальные ($\sim 10$--$45$~млн лет после Большого взрыва). Открытыми остаются точная эпоха появления первого репликатора и вклад дозвёздной/дозёмной химии в суммарную дисперсию.

\item \textbf{Скорость пребиотических шагов.} Лабораторные репликаторы (шаблонная лигация фон Кидровского, пептидные репликаторы, системная химия) реплицируются преимущественно \emph{параболически}, а не экспоненциально, из-за ингибирования продуктом~\cite{Adamski2020}~--- что согласуется с тем, что добиотические шаги медленнее и шумнее; но чистого универсального времени шага для этой стадии данные пока не дают.

\item \textbf{Многофазная структура усложнения и чувствительность $\sigma$.} Часы Шарова~\cite{Sharov2013} и индекс сборки~\cite{Marshall2021,Cronin2023} указывают на $\geq 2$ режима роста сложности: абиотический ($\sim 0{,}6$~Gyr$^{-1}$) и биологический ($\sim 2{,}7$~Gyr$^{-1}$). Единый усреднённый такт $\mu$~--- упрощение. Грубая оценка чувствительности: по долям времени (\S~\ref{sec:mu_from_sigma}) пребиотическая фаза даёт $\sim 65\%$ шагов, ранняя планетная $\sim 9\%$, земная микробная $\sim 22\%$; если собственный такт пребиотики отличается от биологического, $\sigma = \sqrt{\sum_i N_i\mu_i^2}$ сдвигается, но при отличии в пределах заявленной порядковой неопределённости (фактор $\sim 1{,}5$) выводы устойчивы. Отдельно: гауссово приближение хвоста в аргументе тишины (\S~\ref{sec:silence_sigma}) при сильных негауссовых уклонениях способно сместить границу $n$ на порядки~--- это более существенный источник неопределённости, чем многофазность.

\item \textbf{Унифицированные часы сложности.}\label{sec:unified_clocks} Функциональный геном (Шаров), индекс сборки MA и плотность потока свободной энергии $\Phi_m$ (Чейссон) описывают один критический путь с информационной и энергетической сторон (\S~\ref{sec:complexity_clocks}), но согласование их темпов по фазам (абиотика, биология, техносфера) нетривиально. Перспективная задача~--- построить согласованные трёхкомпонентные «часы сложности» от Большого взрыва до цивилизации: явный перевод между $\log N$, $\log \mathrm{MA}$ и $\log \Phi_m$, калибровка пофазных темпов и проверка, насколько такая шкала улучшает однофазное приближение $\mu \approx \mathrm{const}$. Настоящая работа сознательно использует упрощение; полная унификация~--- задел на будущие исследования.

\item \textbf{Носители пребиотической химии.} Где накапливается сложность к порогу C*~--- пыль (IDPs), кометы, метеориты или поверхность планеты~--- окончательно не установлено. Мы склоняемся к мелкой пыли (\S~\ref{sec:idp_dust}: её много, доставляет органику без нагрева~\cite{Anders1989,ChybaSagan1992,Flynn2004}, ничтожный энергобюджет реализуем на отдельной пылинке), но разные носители дают разные оценки числа независимых событий накопления и разброса времён достижения C*.
\end{itemize}

\bibliographystyle{unsrt}
\bibliography{refs}

\end{document}
